%==============================================================================
%  OpenCFD-EC 1.16a 程序说明（多块/多区域扩展版）
%  编译：  xelatex main.tex   （连续两次以生成目录与交叉引用）
%  或：    ./build.sh
%  注意：  中文用 ctex + Fandol 字体（xeCJK），务必用 xelatex/lualatex，勿用 pdflatex
%  维护：  程序新增功能后，请更新对应章节并在 92_changelog.tex 追加一条记录
%==============================================================================
\documentclass[11pt,a4paper,fontset=fandol]{ctexart}
\usepackage[margin=2.4cm]{geometry}
\usepackage{amsmath,amssymb}
\usepackage{booktabs}
\usepackage{longtable}
\usepackage{array}
\usepackage{listings}
\usepackage{xcolor}
\usepackage{enumitem}
\usepackage{caption}
\usepackage{fancyhdr}
\usepackage[hidelinks,bookmarksnumbered]{hyperref}

\pagestyle{fancy}\fancyhf{}
\fancyhead[L]{\small OpenCFD-EC 1.16a 程序说明}
\fancyhead[R]{\small 第\,\thepage\,页}
\renewcommand{\headrulewidth}{0.4pt}
\setlist{nosep,leftmargin=1.7em}
\captionsetup{font=small}
\newcommand{\code}[1]{\texttt{#1}}

\lstset{
  basicstyle=\ttfamily\footnotesize,
  breaklines=true,
  columns=flexible,
  keepspaces=true,
  showstringspaces=false,
  frame=single,
  framerule=0.3pt,
  rulecolor=\color{gray!60},
  backgroundcolor=\color{gray!5},
  xleftmargin=0.8em,
  framexleftmargin=0.8em,
  captionpos=b
}

\title{\bfseries OpenCFD-EC 1.16a 程序说明\\[0.4em]
 \large 多块/多区域扩展版\\[0.2em]
 \normalsize 固体导热 · 可压高速 · 不可压低速 · 多孔介质 · 跨区域耦合}
\author{基于 Li Xinliang（中科院力学所 LHD）OpenCFD-EC 1.16a 的可压缩密度基基线扩展}
\date{版本 1.0 \quad 2026-09-25}

\begin{document}
\maketitle

\begin{abstract}
\noindent
本手册面向 OpenCFD-EC 1.16a 的\textbf{多块/多区域扩展版}，系统说明：
\begin{enumerate}
  \item 程序的块类型体系（可压流体 / 固体 / 低速不可压 / 多孔介质）与跨区域接口码；
  \item 各求解器的物理模型、离散方法、\textbf{控制参数含义、默认值、选取建议与典型标定经验}；
  \item \code{control.ec}（7 组 namelist）及 \code{bc3d.inp} / \code{bc3d\_interface.inp} /
        \code{material.in} / \code{porous.inp} / \code{solid\_bc.inp} 的写法；
  \item 输出与重启（\code{flow3d.dat} / VTK / \code{flow3d\_node.dat} / \code{field\_restart.dat}）；
  \item 以 \code{cases/fluid\_solid} 的\textbf{三个耦合算例}（码 11 / 12 / 19）为例，
        逐一给出算例设置、逐参数解释、运行命令、时间预算与常见问题。
\end{enumerate}
\textbf{维护约定}：程序新增或修改功能后，必须同步更新本手册对应章节，并在
\hyperref[app:changelog]{附录 C“更新记录”}追加一条。各章的“依据文件”给出了可核对的源码位置。
\end{abstract}

\tableofcontents
\clearpage

\section{概述}

\subsection{程序谱系与目标}
本程序以李新亮（中科院力学所 LHD）的可压缩密度基求解器 \textbf{OpenCFD-EC 1.16a}
为基线，在其上扩展\textbf{多块、多区域混合求解}能力，用于流固共轭传热（CHT）、
高低温/高低速混接、发汗冷却（多孔介质）、流多孔耦合等工程问题：

\begin{itemize}
  \item \textbf{保留基线全部能力}：可压 N-S（MUSCL/WENO、Roe/HLLC/Van Leer、SA/SST 等）、
        多块、多网格、MPI/OpenMP、涡轮机内流、混合 FV/FD（SEC）。
  \item \textbf{新增三类块求解器}：固体导热、不可压低速（SIMPLE/SIMPLEC/人工压缩 AC）、
        多孔介质（Darcy--Brinkman--Forchheimer + LTNE 双温度）。
  \item \textbf{新增跨区域耦合框架}：显式接口码（11--19）+ 逐时间步同时耦合 + 分段交错耦合，
        并支持耦合状态随重启文件保存/恢复。
\end{itemize}

\subsection{块类型}
主程序逐块分派求解器（\code{src/sub\_multi\_region.f90: solver\_one\_block}）。
块类型由 \code{material.in} 的首列给出：

\begin{table}[htbp]\centering\small
\caption{块类型码（\code{src/sub\_modules.f90}）}\label{tab:blocktype}
\begin{tabular}{clll}
\toprule
码 & 名称 & 求解器 & 状态 \\
\midrule
0 & \code{BLOCK\_FLUID}    & 可压密度基 N-S（RANS 框架） & 基线继承，完整 \\
1 & \code{BLOCK\_SOLID}    & 固体导热 FVM（$T_w$/$Q_w$ 热边界） & 本仓库新增 \\
2 & \code{BLOCK\_LOWSPEED} & 不可压低速（SIMPLE/SIMPLEC/AC + Rhie--Chow） & 本仓库新增 \\
3 & \code{BLOCK\_POROUS}   & 多孔介质（DBF + LTNE 双温度） & 本仓库新增 \\
\bottomrule
\end{tabular}
\end{table}

\subsection{跨区域接口码}
接口码写在 \code{bc3d.inp} 中（标识“块的类型对”，\textbf{不含方向}；方向由几何配对的
连接描述符 \code{L1..L3} 决定）：

\begin{table}[htbp]\centering\small
\caption{接口码与实现状态}\label{tab:ifccode}
\begin{tabular}{cll}
\toprule
码 & 含义 & 实现状态 \\
\midrule
11 & 可压流体 $\leftrightarrow$ 固体（CHT） & 已实现（\code{cases/fluid\_solid}）\\
12 & 可压流体 $\leftrightarrow$ 低速       & 已实现（\code{cases/high\_low\_fluid}）\\
13 & 低速 $\leftrightarrow$ 固体（CHT）    & 已实现（\code{cases/bl\_cht}）\\
14 & 固固 & 已实现（缓冲交换）\\
15 & 流流（同类可压） & 已实现（缓冲交换）\\
16 & 低速 $\leftrightarrow$ 多孔           & 已实现 + 定量验证（\code{cases/beavers\_joseph}）\\
17 & 固 $\leftrightarrow$ 多孔             & \textbf{未实现} \\
18 & 多孔 $\leftrightarrow$ 多孔           & \textbf{未实现} \\
19 & 可压 $\leftrightarrow$ 多孔（发汗/热耦合） & 已实现（\code{cases/porous\_fluid\_phaseB} 等）\\
\bottomrule
\end{tabular}
\end{table}

除显式接口码外，也支持\textbf{链接式写法}（物理壁面码 + 相邻块号 \code{nb1}/面号 \code{face1}）：
读入后会自动“提升”为规范接口码（\code{src/sub\_IO.f90} 内 \code{is\_interface\_bc} 判定，
\code{bc<0} 或 $11\le bc\le 19$ 视为界面）。

\subsection{物理边界码}
\begin{table}[htbp]\centering\small
\caption{物理边界码（\code{src/sub\_modules.f90}）}\label{tab:physbc}
\begin{tabular}{cll}
\toprule
码 & 名称 & 说明 \\
\midrule
2 & \code{BC\_Wall}        & 无滑移壁（可压侧 $T_{wall}<0$ 表示绝热）；低速侧 $j+$ 面可由 \code{LS\_U\_lid} 驱动 \\
3 & \code{BC\_Symmetry}    & 对称面（法向速度镜像、切向与标量零梯度）\\
4 & \code{BC\_Farfield}    & 远场/无反射（低速侧按出口处理）\\
5 & \code{BC\_Inflow}      & 入口（可压按来流；低速按 \code{LS\_Inlet\_Type}）\\
6 & \code{BC\_Outflow}     & 出口（可压 $p_{outlet}$ 或外插；低速按给定静压 $+$ 零梯度）\\
21 & \code{BC\_LS\_Inlet}  & 低速专用入口（与码 5 等价处理）\\
22 & \code{BC\_LS\_Outlet} & 低速专用出口（与码 6 等价处理）\\
201 & \code{BC\_Wall\_Turbo} & 涡轮机壁面（旋转壁）\\
401 & \code{BC\_Extrapolate} & 外插 \\
501 & \code{BC\_Periodic}    & 周期性（配合 \code{Periodic\_dX/dY/dZ}）\\
$-1$ & \code{BC\_Inner}      & 块间内部界面（由 \code{bc3d\_interface.inp} 给出配对）\\
$-2/-3$ & \code{BC\_PeriodicL/R} & 周期界面左/右 \\
\bottomrule
\end{tabular}
\end{table}

\subsection{面编号约定（全局统一，务必牢记）}
\begin{center}
\code{1=i-, 2=j-, 3=k-, 4=i+, 5=j+, 6=k+}
\end{center}
\code{bc3d.inp} 每行给出 $i_1,i_2,j_1,j_2,k_1,k_2$（节点编号范围）与面码；
\code{bc3d\_interface.inp} 同时给出连接描述符 \code{L1,L2,L3}：本块第 $k$ 维与邻块第 $|L_k|$ 维相连，
符号表示正/反向连接（\code{src/sub\_convert\_inp.f90: Convert\_bc}）。

\subsection{单位与无量纲化（最易出错）}
\begin{table}[htbp]\centering\small
\caption{各块的变量与单位约定}\label{tab:units}
\begin{tabular}{p{2.6cm}p{5.4cm}p{6.6cm}}
\toprule
块类型 & 存什么（\code{U}） & 单位/无量纲化 \\
\midrule
可压流体 & \code{U=(rho, rho*u, rho*v, rho*w, E)}；$T^*=T/T_\infty$ & 无量纲：$\rho_\infty$、$U_\infty=Ma\,a_\infty$、$p^*=p/(\rho_\infty U_\infty^2)$；网格坐标乘 \code{Lscale} 得 SI \\
固体     & 温度在 \code{B\%Ts}（K） & \textbf{SI}（K；$k$ 用 W/(m$\cdot$K)，$\rho$、$C_p$ 用 SI）\\
低速     & \code{U(1)=rho=LS\_rho}（常密度）、\code{U(2..4)=速度 [m/s]}、\code{U(5)=T [K]}；压力在 \code{B\%p} & \textbf{SI}（注意：低速块存\textbf{速度}而非动量）\\
多孔     & 同低速（流体温度 $T_f$ 在 \code{U(5)}），骨架温度在 \code{B\%Ts} & \textbf{SI}；\code{LS\_rho} 为流体密度参考 \\
\bottomrule
\end{tabular}
\end{table}

可压参考量（\code{src/sub\_read\_parameter.f90: set\_const\_para}）：
$a_{ref}=\sqrt{\gamma R T_\infty}$，$U_{ref}=Ma\,a_{ref}$，
$\rho_{ref}=\dfrac{Re\,\mu_\infty(T_\infty)}{Ma\,a_{ref}\,\texttt{Lscale}}$，
$p_{ref}=\rho_{ref}U_{ref}^2$。
注意 \textbf{$Re$ 以 \code{Lscale} 为参考长度}：\code{Lscale}=0.001（毫米网格）时 $Re=5000$
对应板长雷诺数约 $8.3\times10^{5}$。

\subsection{目录结构}
\begin{lstlisting}
src/                 主程序 + 各子模块（含 makefile），产物 opencfd-ec1.16a.out
cases/               算例（网格/边界/物性/control.ec）
docs/                文档；程序说明源码在 docs/程序说明/，PDF 在 docs/
util/                后处理/网格工具（readflow3d、check_flow3d_node.py 等）
memory-bank/         开发记忆库（projectbrief/activeContext/progress/worklog/constraints）
\end{lstlisting}

\subsection{最小上手流程}
\begin{enumerate}
  \item \code{cd src \&\& make}（默认 \code{mpif90}，产物 \code{opencfd-ec1.16a.out}）；
  \item 复制 \code{control.ec.template} 到算例目录，按需保留/修改各 namelist 组；
  \item 准备 \code{Mesh3d.x}、\code{bc3d.inp}、\code{bc3d\_interface.inp}、\code{material.in}
        （多孔/固体再加 \code{porous.inp}、\code{solid\_bc.inp}）；
  \item \code{cd cases/<case> \&\& mpirun -np 1 ./opencfd-ec1.16a.out}
        （交错耦合模式必须单进程，详见第 \ref{sec:build} 章）。
\end{enumerate}

\noindent\textbf{依据文件}：\code{src/sub\_modules.f90}、\code{src/sub\_multi\_region.f90}、
\code{src/sub\_convert\_inp.f90}、\code{src/sub\_IO.f90}、\code{docs/control.ec-说明.md}。

\section{输入文件与算例结构}

\subsection{算例目录文件清单}
\begin{table}[htbp]\centering\small
\caption{算例目录中的文件}\label{tab:casefiles}
\begin{tabular}{p{4.2cm}p{1.3cm}p{9.5cm}}
\toprule
文件 & 必需 & 说明 \\
\midrule
\code{control.ec} & 必需 & 全部求解器控制参数（7 组 namelist，见第 \ref{sec:params} 节与附录 A）\\
\code{Mesh3d.x} / \code{Mesh3d.dat} & 必需 & 网格；由 \code{\$flow\_ec: Mesh\_File\_Format} 选择格式（\code{2}/\code{0}）\\
\code{bc3d.inp} & 必需 & 逐块逐面物理边界（可含接口码 11--19）\\
\code{bc3d\_interface.inp} & 多块必需 & 块间对接关系（配对 + 连接描述符）\\
\code{material.in} & 多区域必需 & 块类型表；固体/多孔块另给 $(\rho, C_p, k)$ \\
\code{solid\_bc.inp} & 有固体/多孔时 & 逐面热边界：等温 $T_w$ 或热流 $Q_w$ \\
\code{porous.inp} & 有多孔块时 & 多孔块 $\varepsilon,\ d_p,\ h_v$ \\
\code{partation.dat} & 可选 & 进程分区；交错耦合模式\textbf{不要}放此文件 \\
\code{field\_restart.dat} & 自动生成 & 重启文件；存在即自动续算（\code{Iflag\_restart=0}）\\
\code{bc3d.inc} / \code{bc3d\_interface.inc} & 自动生成 & 由 \code{.inp} 转换而来（含 \code{L1..L3}）\\
\bottomrule
\end{tabular}
\end{table}

\subsection{\code{control.ec}：7 组 namelist}
\label{sec:params}
历史版本把所有参数放在\textbf{一个} \code{\$control\_ec}（128 个变量）里。现拆成 7 组，
\textbf{每组、每个变量都可省略}（全部有代码默认值），算例只需写与默认值不同的项。

\begin{table}[htbp]\centering\small
\caption{\code{control.ec} 的 7 个组}\label{tab:groups}
\begin{tabular}{lp{5.2cm}cp{5.2cm}}
\toprule
组名 & 用途 & 变量数 & 关键变量（详见附录 A） \\
\midrule
\code{\$freestream\_ec} & 来流/参考态/气体物性 & 24 & \code{Ma, Re, AoA, AoS, T\_inf, Twall, p\_outlet, gamma, PrL, Lscale} \\
\code{\$flow\_ec}       & 时间步/格式/模型/网格/输出/重启 & 52 & \code{t\_end, CFL, dt\_global, Time\_Method, Iflag\_Scheme, Iflag\_Flux, If\_viscous, Mesh\_File\_Format, Kstep\_save, Iflag\_restart} \\
\code{\$lowspeed\_ec}   & 低速（不可压）块 & 21 & \code{LS\_rho, LS\_mu, LS\_k, LS\_Cp, LS\_Inlet\_Type, LS\_U\_in, LS\_P\_out, LS\_Max\_Iter, LS\_Algorithm} \\
\code{\$ac\_ec}         & 人工压缩(AC)求解器（低速与多孔共用） & 15 & \code{AC\_Max\_Iter, AC\_beta, AC\_CFL, AC\_CFLv, AC\_Tol, AC\_Flux, AC\_Recon, AC\_Limiter} \\
\code{\$solid\_ec}      & 固体导热 GS 控制 & 4 & \code{Solid\_GS\_Omega, Solid\_Max\_Iter, Solid\_Min\_Iter, Solid\_Tol} \\
\code{\$porous\_ec}     & 多孔块 & 8 & \code{Porous\_T\_ref, Porous\_alpha\_Ts, Porous\_Max\_Iter, Porous\_Tol, Porous\_U/V/W\_in, Porous\_T\_in} \\
\code{\$couple\_ec}     & 跨区域耦合调度 & 12 & \code{Iflag\_Couple\_Scheme, Kstep\_Couple\_Comp, Niter\_Couple\_Outer, Tol\_Couple\_Tw, Iflag\_Couple\_Restart} \\
\bottomrule
\end{tabular}
\end{table}

\paragraph{兼容与读取规则}
\begin{enumerate}
  \item 旧组 \code{\$control\_ec}（128 变量 $+$ \code{Solid\_*}）\textbf{完全保留}，老算例零改动即可运行；
  \item 启动时先\textbf{扫描}文件判定存在哪些组（只有“首个非空字符为 \code{\$}/\code{\&}”的行算组头；
        \code{!} 之后为注释；大小写不敏感），再按 旧组 $\to$ \code{freestream} $\to$ \code{flow} $\to$
        \code{lowspeed} $\to$ \code{ac} $\to$ \code{solid} $\to$ \code{porous} $\to$ \code{couple} 顺序读取；
        \textbf{分组值覆盖旧组值}（支持逐组渐进迁移）；
  \item 某组存在但解析失败（变量名拼错、缺 \code{\$end}）$\Rightarrow$ \textbf{报错并 \code{stop 1}}，
        不会静默回落默认值；
  \item 启动日志与 \code{output\_para.out} 末尾会回显各组来源（\code{T}=来自文件，\code{F}=默认）
        与固体 GS 四参数，便于回归比对。
\end{enumerate}

\paragraph{最小示例（3 组即可跑一个流固耦合算例）}
\begin{lstlisting}
$freestream_ec
  Ma=3.0d0, Re=5000.d0, T_inf=800.d0, Lscale=0.001d0
$end
$flow_ec
  t_end=1.0d6, Kstep_save=500, Kstep_show=50,
  dt_global=1.0d0, Mesh_File_Format=2, Iflag_vtk_onefile=1, Iflag_vtk_SI=1,
  Iflag_restart=0, Kstep_restart=0, Iflag_flow_node=1
$end
$couple_ec
  Iflag_Couple_Scheme=1, Niter_Couple_Warm=2, Niter_Couple_Outer=12, Tol_Couple_Tw=1.0d-2
$end
\end{lstlisting}
模板见仓库根目录 \code{control.ec.template}（含全部变量、默认值与中文注释）。



\subsection{网格文件}
\begin{itemize}
  \item \code{Mesh\_File\_Format=0}：\code{Mesh3d.dat}（unformatted）；
  \item \code{Mesh\_File\_Format=1}：\code{Mesh3d.dat}（formatted），可用文本工具查看；
  \item \code{Mesh\_File\_Format=2}：\code{Mesh3d.x}（PLOT3D 风格，\textbf{x、y、z 三个变量逐块一条记录}）。
\end{itemize}
多块网格在同一文件中按块依次存放；块序即输出块序（VTK、\code{flow3d\_node.dat}、重启文件同序）。
网格坐标用\textbf{网格单位}（本仓库算例为 mm，配合 \code{Lscale=0.001}）。

\subsection{\code{bc3d.inp}：物理边界}
格式（块循环）：
\begin{lstlisting}
nblock
  ni nj nk           (第 1 块维数)
  blk-1              (块名)
  nface              (该块面数)
    i1 i2 j1 j2 k1 k2 bc      (重复 nface 次)
  ni nj nk           (第 2 块维数) ... 依此类推
\end{lstlisting}
\code{i1..k2} 为\textbf{节点}编号范围；\code{bc} 见表 \ref{tab:physbc} 与 \ref{tab:ifccode}。
例如 case1 固体块（Block 1，$67\times101\times2$）的 $j-$ 面写成接口码 11：
\begin{lstlisting}
       1       67        1        1        1        2        11
\end{lstlisting}
而可压块（Block 2）的 $i+$ 面（$j=68..134$ 段）也写 11——两块互相配对。

\subsection{\code{bc3d\_interface.inp}：块间对接}
结构与 \code{bc3d.inp} 相同，但用\textbf{链接描述符}表达“哪个面连到哪个块”：
\begin{lstlisting}
i1 i2 j1 j2 k1 k2   bc (= -1 内部界面, -2/-3 周期)
nb1   : 邻块号（跨行续写时前一行可写 -1，下一行给出）
face1 : 邻块上的对应面号
L1 L2 L3 : 连接描述符（本块第 k 维 <-> 邻块第 |Lk| 维，符号=正/反向）
\end{lstlisting}
示例（case1 固体块 $j-$ 面连到块 2 的 $i+$ 面、切向反向）：
\begin{lstlisting}
        1       67        1        1       -1       -2       -1
        2        4     -2    1        3
\end{lstlisting}
转换后 \code{bc3d\_interface.inc} 给出每个面的 \code{L1,L2,L3}；耦合例程\textbf{完全依赖 L 表}
做指标映射，因此新拓扑一般无需改物理公式，只需支持新的 $(face,\,face1)$ 组合（见第 \ref{sec:coupling} 章）。

\subsection{\code{material.in}：块类型与物性}
\begin{lstlisting}
nblock
  blk1_type  blk2_type  ...        (一行给出全部块类型码 0/1/2/3)
  rho1 Cp1 k1                      (块 1 物性; 固体/多孔块用 SI)
  rho2 Cp2 k2 ...                  (逐块一行; 流体/低速块可写 0 0 0)
\end{lstlisting}
case1 实例（块 1 = 不锈钢固体、块 2 = 空气）：
\begin{lstlisting}
2
1 0
7900.0 500.0 16.3
0.0 0.0 0.0
\end{lstlisting}
多孔块的“骨架”物性同样写在此文件（类型码 3），孔尺度/换热系数写 \code{porous.inp}。

\subsection{\code{porous.inp}：多孔块参数}
\begin{lstlisting}
nblock
  m  eps  dp[m]  hv[W/(m^3 K)]      (块号, 孔隙率, 颗粒/孔径尺度, 体积换热系数)
\end{lstlisting}
case3 实例：块 1、$\varepsilon=0.30$、$d_p=1$ mm、$h_v=2\times10^{6}$ W/(m$^{3}$K)。
渗透率按 Ergun 关系由 $d_p,\varepsilon$ 计算（第 \ref{sec:porous} 章）；\code{hv}=0 退化为
局部热平衡（$T_s\equiv T_f$）。

\subsection{\code{solid\_bc.inp}：固体/多孔骨架热边界}
\begin{lstlisting}
nSolidBlocks
  m                 (块号)
  nf                (该块给热边界的面数)
  face  Tw[K]  Qw[W/m^2]            (重复 nf 次; 可选再跟 htc/Tinf 列)
\end{lstlisting}
约定：\textbf{$T_w>0$ 表示等温壁（K）；$T_w<0$ 表示给定热流 $Q_w$（W/m$^{2}$）}。
case1 实例（固体块 1 的 $j+$ 面等温 300 K）：
\begin{lstlisting}
1
1
1
5 300.0 0.0
\end{lstlisting}

\noindent\textbf{依据文件}：\code{src/sub\_read\_parameter.f90}、\code{src/sub\_convert\_inp.f90}、
\code{src/sub\_IO.f90}、\code{cases/fluid\_solid/*}、\code{control.ec.template}。
%===EOF===

\section{可压高速求解器（\code{BLOCK\_FLUID}）}
\label{sec:compressible}

\subsection{变量与无量纲化}
块类型 0，\textbf{密度基可压 N-S}（RANS 框架）。状态量
$U=(\rho,\ \rho u,\ \rho v,\ \rho w,\ E)$，无量纲参考量见 \S 1.6：
$\rho_\infty$、$U_\infty=Ma\,a_\infty$、$T_\infty=\texttt{T\_inf}$、$p^*=p/(\rho_\infty U_\infty^2)$。
粘性系数由 Sutherland 律 + \code{T\_inf} 定标；层流/湍流 Prandtl 数为 \code{PrL}/\code{PrT}。
网格坐标乘 \code{Lscale} 得 SI（因此 $Re$ 的参考长度是 \code{Lscale}）。

\subsection{时间推进}
\begin{table}[htbp]\centering\small
\caption{时间推进控制参数（组 \code{\$flow\_ec}）}\label{tab:time}
\begin{tabular}{llp{8.4cm}}
\toprule
参数 & 默认值 & 说明/选取建议 \\
\midrule
\code{Time\_Method} & 0 & 0 $=$ LU-SGS 隐式（稳态首选）；1 $=$ 一步 Euler；3 $=$ RK3；
      $-1$ $=$ 双时间步（内迭代）\\
\code{CFL} & 1.0 & 局部时间步 CFL；定常问题可取 2--5（隐式较鲁棒）\\
\code{dt\_global} & 0.01 & 全局时间步（无量纲）；非定常/交错耦合里常用 1.0 与 \code{tt} 同步计数\\
\code{Iflag\_local\_dt} & 1 & 1 $=$ 局部时间步（加速收敛，非定常需置 0）\\
\code{dtmax}/\code{dtmin} & 10/1e-9 & 时间步上下限\\
\code{w\_LU} & 1.0 & LU-SGS 松弛因子 \\
\code{Step\_Inner\_Limit}/\code{Res\_Inner\_Limit} & 20 / 1e-10 & 双时间步内迭代上限/内残差判据 \\
\code{If\_Residual\_smoothing} & 0 & 残差光滑开关 \\
\code{If\_dtime\_mesh} & 1 & 网格质量差时自动减小时间步 \\
\bottomrule
\end{tabular}
\end{table}

\subsection{空间离散}
\begin{table}[htbp]\centering\small
\caption{空间格式参数}\label{tab:scheme}
\begin{tabular}{llp{8.2cm}}
\toprule
参数 & 默认值 & 可选值与建议 \\
\midrule
\code{Iflag\_Scheme} & 5 & 0 一阶迎风；1 NND2；2 UD3；3 MUSCL2U；4 MUSCL2C；5 MUSCL3；6 OMUSCL2；
      7 WENO5；8 UD5；9 WENO7。\textbf{激波/高超声速建议 5 或 7}；低速内流可用 3 \\
\code{Iflag\_Flux} & 5 & 1 Steger--Warming；2 HLL；3 HLLC；4 Roe；5 Van Leer；6 AUSM。\\
\code{IFlag\_Reconstruction} & 0 & 0 原始变量；1 守恒变量；2 特征变量重构 \\
\code{Bound\_Scheme} & MUSCL2C & 边界插值格式（\code{Scheme\_MUSCL2C}）\\
\code{IF\_Scheme\_Positivity} & 1 & 正值保持（1 开）；发散/负密度时优先保留开启 \\
\code{Ldmin,Ldmax} & 1e-6, 1000 & 限制器参数（密度） \\
\code{Lpmin,Lpmax} & 1e-6, 1000 & 压力限制 \\
\code{Lumax,LSAmax} & 1000, 1000 & 速度/湍流量限制 \\
\bottomrule
\end{tabular}
\end{table}

\subsection{湍流模型与壁距}
\begin{itemize}
  \item \code{Iflag\_turbulence\_model}：0 无；1 BL（代数）；2 SA / New-SA；3 SST（$k$--$\omega$）。
        默认 0。若需要热流（壁面 $q_w$）\textbf{必须有粘性}：\code{If\_viscous=1}。
  \item \code{MUT\_MAX}：涡粘上限（$<0$ 不限制），用于极端算例防发散；
        \code{Kt\_inf}/\code{Wt\_inf} 为 SST 初值；\code{CP1\_NSA}/\code{CP2\_NSA} 为 New-SA 常数。
  \item 壁距由程序自动计算（\code{wall\_dist.dat}，缺失时自动生成）。
\end{itemize}

\subsection{来流/参考态与气动力}
\begin{table}[htbp]\centering\small
\caption{来流与参考量（组 \code{\$freestream\_ec}）}\label{tab:freestream}
\begin{tabular}{llp{8.2cm}}
\toprule
参数 & 默认值 & 说明 \\
\midrule
\code{Ma} & 1.0 & 来流马赫数（可压块）\\
\code{Re} & 1000 & 雷诺数，\textbf{参考长度为 \code{Lscale}} \\
\code{T\_inf} & 288.15 & 来流/参考温度 [K]（也是粘性律参考温度、固体块 \code{Ts} 初值）\\
\code{AoA}/\code{AoS} & 0 & 攻角/侧滑角 [deg] \\
\code{Twall} & $-1$ & 可压侧物理壁温 [K]；$<0$ 表示绝热壁 \\
\code{p\_outlet} & $-1$ & 出口压力（无量纲）；$<0$ 表示外插 \\
\code{gamma},\code{PrL},\code{PrT} & 1.4, 0.7, 0.9 & 比热比、层流/湍流 Prandtl 数 \\
\code{Lscale} & 1.0 & 长度尺度（1=m，0.001=mm）\\
\code{Ref\_S}/\code{Ref\_L}/\code{Centroid} & 1/1/0 & 气动力/力矩参考面积、长度、矩心 \\
\code{Cood\_Y\_UP} & 1 & 坐标系 $y$ 轴朝上标志 \\
涡轮机族 & --- & \code{IF\_TurboMachinary, Ref\_medium\_usrdef, Turbo\_P0/T0/L0/w, Turbo\_Periodic\_seta}（内流模式，见原版手册）\\
\bottomrule
\end{tabular}
\end{table}

\subsection{网格、并行与输出节奏}
\begin{table}[htbp]\centering\small
\caption{网格/并行/输出参数（组 \code{\$flow\_ec}）}\label{tab:flowio}
\begin{tabular}{llp{8.0cm}}
\toprule
参数 & 默认值 & 说明 \\
\midrule
\code{Mesh\_File\_Format} & 0 & 0 unformatted \code{Mesh3d.dat}；1 formatted；\textbf{2 \code{Mesh3d.x}（本仓库算例）}\\
\code{Num\_Mesh}/\code{Pre\_Step\_Mesh} & 1/0 & 多网格层数与各层预推进步数 \\
\code{NUM\_THREADS} & 1 & OpenMP 线程数（需 \code{-openmp} 编译）\\
\code{IFLAG\_LIMIT\_FLOW} & 1 & 流场限制（防止负密度/负压）；算例常置 0 以复现基线结果 \\
\code{Iflag\_init} & 0 & 初场方式（0 均匀场；1 由 \code{flow3d.dat} 读入；$-1$ 置零）\\
\code{Kstep\_save} & 1000 & 输出间隔（\code{flow3d.dat}/VTK/\code{Ts\_block}/\code{flow3d\_node}/\code{field\_restart} 同节奏）\\
\code{Kstep\_show} & 1 & 残差/气动力打印间隔 \\
\code{Kstep\_average} & 0 & 时间平均输出间隔（0 关闭）\\
\code{Kstep\_smooth}/\code{Kstep\_init\_smooth} & $-1$/0 & 流场光滑间隔（$<0$ 关闭） \\
\code{Iflag\_savefile} & 0 & 0 关闭、1 开启 \code{flow3d.dat} \\
\code{Iflag\_vtk\_onefile} & 0 & 0 每块一个 \code{flow3d\_block\_*.vtk}；1 合并为 \code{flow3d.vtk}（约 11 MB 级）\\
\code{Iflag\_vtk\_SI} & 0 & 1 $=$ VTK 输出换 SI 单位（含低速/固体块直接输出 SI）\\
\code{Iflag\_bc\_check} & 1 & 0 旧行为；1 自动生成/校验 \code{bc3d.inp}；2 严格报错 \\
\code{Periodic\_dX/dY/dZ} & 0 & 周期平移量 \\
\bottomrule
\end{tabular}
\end{table}

\subsection{参数选取建议（经验值）}
\begin{itemize}
  \item \textbf{超声速/含激波}：\code{Iflag\_Scheme=5}（MUSCL3）或 7（WENO5）、
        \code{Iflag\_Flux=5}（Van Leer）或 4（Roe），\code{CFL=1--2}，\code{Time\_Method=0}。
  \item \textbf{粘性热流/CHT}：务必 \code{If\_viscous=1}；壁温由 \code{Twall}（$<0$ 绝热）或
        跨区域耦合给出；湍流可选 2（SA）或 3（SST）。
  \item \textbf{首次调试}：先用较小 \code{Kstep\_save}/\code{Kstep\_show} 看残差，再放大；
        \code{IF\_Debug=1} 可打印接口配对诊断（\code{dbg\_dump\_interfaces}）。
  \item \textbf{交错耦合模式}：不要用 \code{t\_end} 控时（由 \code{Niter\_Couple\_Outer} 控制），
        \code{t\_end} 给大值即可（详见第 \ref{sec:coupling} 章）。
\end{itemize}

\noindent\textbf{依据文件}：\code{src/sub\_read\_parameter.f90}、\code{src/sub\_flux\_split.f90}、
\code{src/sub\_time\_advance.f90}、\code{src/sub\_LU\_SGS.f90}、\code{docs/程序能力与算例考核总结.md}。
%===EOF===

\section{低速不可压求解器（\code{BLOCK\_LOWSPEED}）}
\label{sec:lowspeed}

\subsection{变量约定（易错，务必先读）}
块类型 2。原始变量\textbf{格心存储}：
\begin{center}
\code{U(1) = rho = LS\_rho}（常密度）、\ \code{U(2..4) = 速度分量 [m/s]}、\ \code{U(5) = T [K]}；
压力在 \code{B\%p}（含鬼点层）。
\end{center}
\textbf{注意：低速块存的是速度而不是动量}（\code{rho*u}）。低速/多孔块的物性与结果全部为 \textbf{SI}：
长度用 \code{mesh\,[grid]} $\times$ \code{Lscale}[m]，速度 m/s，温度 K，压力 Pa。

\subsection{三种压力--速度耦合算法（\code{LS\_Algorithm}）}
\begin{table}[htbp]\centering\small
\caption{\code{LS\_Algorithm} 选择}\label{tab:lsalg}
\begin{tabular}{clp{8.6cm}}
\toprule
值 & 算法 & 说明/适用 \\
\midrule
1 & SIMPLE  & 经典稳态压力修正；对角占优好、网格均匀的算例（如 \code{lidcavity}）\\
2 & SIMPLEC & 对动量方程邻居项做一致化处理，收敛略快；强耦合/细长通道可试 \\
3 & \textbf{AC-FV} & 人工压缩（虚拟压缩）伪时间推进 $+$ LU-SGS（\code{src/sub\_lowspeed\_ac.f90}）。
    对\textbf{强耦合、长通道、低速流固/流多孔}更鲁棒，是 \code{cases/fluid\_solid} 的\textbf{必需}选择\\
\bottomrule
\end{tabular}
\end{table}
SIMPLE/SIMPLEC 用 \textbf{Rhie--Chow} 插值抑制棋盘压力振荡，动量/压力/温度分别有欠松弛
\code{LS\_alpha\_u}/\code{LS\_alpha\_p}/\code{LS\_alpha\_T}；
每次调用求解器内部迭代 \code{LS\_Max\_Iter} 次或压力修正残差 $<$ \code{LS\_Tol}。
对流项格式由 \code{LS\_Scheme} 选择：1 一阶迎风（最稳）、2 二阶迎风、3 MUSCL(Van Leer)。

\subsection{低速参数表}
\begin{table}[htbp]\centering\small
\caption{组 \code{\$lowspeed\_ec}（全部 SI）}\label{tab:lsparams}
\begin{tabular}{llp{8.0cm}}
\toprule
参数 & 默认值 & 说明/选取建议 \\
\midrule
\code{LS\_rho} & 1.2 & 密度 [kg/m$^3$]（\textbf{常密度}）\\
\code{LS\_mu}  & 1.8e-5 & 动力粘度 [Pa$\cdot$s] \\
\code{LS\_k}   & 0.025 & 导热系数 [W/(m$\cdot$K)]（温度方程/CHT 用）\\
\code{LS\_Cp}  & 1005 & 比热 [J/(kg$\cdot$K)] \\
\code{LS\_T\_ref} & 288.15 & 参考/入口温度 [K]（初场、多孔 \code{Porous\_T\_in} 缺省值）\\
\code{LS\_Inlet\_Type} & 1 & 1 速度入口；2 质量流量入口；3 压力入口（\textbf{见下方限制}）\\
\code{LS\_U\_in},\code{LS\_V\_in},\code{LS\_W\_in} & 1,0,0 & 入口速度分量 [m/s] \\
\code{LS\_Mdot\_in} & 0 & 质量流量 [kg/s]（类型 2）\\
\code{LS\_P\_in} & 101325 & 入口总压 [Pa]（类型 3）\\
\code{LS\_P\_out} & 101325 & 出口静压 [Pa]（出口面：给定静压 $+$ 零梯度外推）\\
\code{LS\_T\_wall} & $-1$ & 壁温 [K]；$<0$ 绝热 \\
\code{LS\_U\_lid} & 0 & 顶盖驱动速度 [m/s]（$j+$ 面）\\
\code{LS\_alpha\_p},\code{LS\_alpha\_u},\code{LS\_alpha\_T} & 0.3,0.7,0.7 & 欠松弛因子（SIMPLE/SIMPLEC）\\
\code{LS\_Max\_Iter},\code{LS\_Tol} & 5000, 1e-8 & 每次调用的内部迭代上限/压力修正残差判据 \\
\code{LS\_Scheme} & 1 & 1 一阶迎风；2 二阶迎风；3 MUSCL \\
\code{LS\_Algorithm} & 1 & 1 SIMPLE；2 SIMPLEC；\textbf{3 AC-FV} \\
\bottomrule
\end{tabular}
\end{table}

\paragraph{入口类型的使用限制（重要）}
物理边界码 5 / 21 均触发入口处理（\code{BC\_Inflow} / \code{BC\_LS\_Inlet}），出口用 6 / 22。
当前实现中 \code{LS\_Inlet\_Type=2}（质量入口）按\textbf{$i$ 面 $+$ 网格单位面积}处理，
用在 $j$ 面上方向/量纲会错。工程做法：把目标质量通量 $G$ [kg/(m$^2$s)] 换算为等效法向速度
\begin{equation}
\code{LS\_V\_in} = \frac{G}{\code{LS\_rho}}
\end{equation}
并用速度入口（类型 1）。例如 case2/case3 中 $G=2$ kg/(m$^2$s)、$\rho=1.177$ kg/m$^3$
$\Rightarrow$ 1.699 m/s（$+y$ 方向）。\textbf{该项已列入待办}（改为按面法向赋速度并明确面积量纲）。

\subsection{AC（人工压缩）求解器参数（\code{LS\_Algorithm=3}）}
伪时间推进：$\dfrac{\partial q}{\partial \tau}+\beta\nabla\!\cdot\!\mathbf{u}=\cdots$，
其中 $q=p/\rho$，$\beta=\code{AC\_beta}\cdot U_{ref}^2$（$U_{ref}$ 由入口速度/顶盖速度/
$\sqrt{|\Delta p|/\rho}$ 依次兜底估计）。默认参数：

\begin{table}[htbp]\centering\small
\caption{组 \code{\$ac\_ec}（低速与多孔共用）}\label{tab:acparams}
\begin{tabular}{llp{8.2cm}}
\toprule
参数 & 默认值 & 说明 \\
\midrule
\code{AC\_Max\_Iter} & 40000 & 每次调用的伪时间迭代上限 \\
\code{AC\_Print} & 500 & 打印间隔 \\
\code{AC\_beta} & 10 & $\beta=\code{AC\_beta}U_{ref}^2$（越大越\"定常\"、越稳但越慢）\\
\code{AC\_CFL} & 2 & 无粘 CFL；\textbf{大 CFL 会数轮后发散，需按算例标定} \\
\code{AC\_CFLv} & 0.5 & 粘性 CFL \\
\code{AC\_Tol} & 1e-7 & 收敛判据（无量纲残差）\\
\code{AC\_w} & 1.0 & LU-SGS 松弛 \\
\code{AC\_Flux} & 1 & 1 Rusanov(LLF)；2 Steger--Warming（3 $=$ AUSM+ 已于 2026-09-17 移除，会告警回落）\\
\code{AC\_Recon} & 1 & 1 二阶 MUSCL(van Leer)；2 WENO5；3 WENO3 \\
\code{AC\_Limiter} & 1 & 1 van Leer；2 minmod \\
\code{AC\_WenoBlend} & 0 & WENO 路径混入的一阶 Rusanov 比例（保正性用）\\
\code{AC\_WallRecon} & 1 & 壁/对称面：切向状态单侧二阶重构 \\
\code{AC\_WallP} & 1 & 壁/对称面：压力二阶线性重构 \\
\code{AC\_MomDiss} & 0 & 动量耗散：0 lumped（默认）；1 分裂混合；2 纯 $|u_n|$（\textbf{不稳定}）\\
\code{AC\_MomFrac} & 0.2 & \code{AC\_MomDiss=1} 时保留的声速比例下限 \\
\bottomrule
\end{tabular}
\end{table}
越界/非法取值由 \code{ac\_check\_controls()} 打印告警并回落到安全默认值（不再静默使用未定义值）。

\paragraph{AC 标定经验（来自已有算例，务必参考）}
\begin{table}[htbp]\centering\small
\caption{AC 参数的算例标定记录}\label{tab:accalib}
\begin{tabular}{lll}
\toprule
算例 & 关键设置 & 结论 \\
\midrule
\code{lidcavity\_ac} & \code{AC\_CFL=20}, \code{AC\_beta=1}, \code{LS\_Algorithm=3} & 稳定，$Re=100$ 与 Ghia 吻合 \\
\code{channel/type*} & \code{AC\_*}\ 默认附近 & 三类入口均与 Poiseuille 解析吻合 \\
\code{high\_low\_fluid\_800K} & \code{AC\_beta=10}, \code{AC\_CFL=20}, \code{AC\_Max\_Iter=20000} & 配 \code{LS\_U\_in=10} 作 $U_{ref}$ $\Rightarrow \beta\approx1000$ m$^2$/s$^2$ \\
\textbf{\code{fluid\_solid/run\_case2}} & \textbf{\code{AC\_CFL=2}, \code{AC\_beta=10}, \code{AC\_Max\_Iter=5000}} & \code{AC\_CFL=20} 会数轮后 NaN；\code{LS\_Algorithm=1}(SIMPLE) \textbf{第 1 步即发散} \\
\code{beavers\_joseph/run\_ac} & \code{AC\_CFL=20}, \code{AC\_beta=1000}, \code{AC\_Max\_Iter=1} & 多块 AC 即时交换接口 ghost 的用法 \\
\bottomrule
\end{tabular}
\end{table}

\subsection{温度方程与共轭传热}
低速块带温度方程（$U(5)=T$，物性 \code{LS\_k}/\code{LS\_Cp}）。与固体块（码 13）或
可压块（码 12）对接时，界面温度/热流由耦合例程按热流平衡给出（见第 \ref{sec:coupling} 章）。
壁面热边界由 \code{LS\_T\_wall}（$<0$ 绝热）或跨区域耦合决定。

\subsection{选参建议}
\begin{itemize}
  \item \textbf{先判来源}：均匀网格 + 单块 $\to$ SIMPLE 即可；长通道、强耦合、多块、低速流固
        $\to$ 直接用 AC-FV（\code{LS\_Algorithm=3}）。
  \item \textbf{AC 稳定三步法}：先 \code{AC\_CFL=2}（保守）跑通 $\to$ 若收敛慢再升到 10--20
        $\to$ 若出现 NaN 立即回退并调小 \code{AC\_beta}（$\beta$ 大则伪时间步小、更稳）。
  \item \textbf{残差判据}：\code{AC\_Tol=1e-6}（粗算）$\sim$ \code{1e-7}（默认）；
        \code{AC\_Max\_Iter} 与多孔段的 \code{Porous\_Chunk\_Iter} 共同决定耦合预算（见 \S\ref{sec:coupling}）。
  \item \textbf{已知限制}：\code{cases/flatplate\_re1e6} 与 \code{high\_low\_fluid} 存在低速网格病态
        （收敛慢/精度损失）；SIMPLE 压力求解器为 Jacobi/Gauss--Seidel 型，
        后续可换 BiCGStab/ILU（见第 \ref{sec:validation} 章待办）。
\end{itemize}

\noindent\textbf{依据文件}：\code{src/sub\_lowspeed.f90}、\code{src/sub\_lowspeed\_ac.f90}、
\code{docs/AC-FV-后续工作.md}、\code{docs/程序能力与算例考核总结.md}。
%===EOF===

\section{多孔介质求解器（\code{BLOCK\_POROUS}）}
\label{sec:porous}

\subsection{物理模型}
块类型 3。不可压 + \textbf{Darcy--Brinkman--Forchheimer (DBF)} 动量方程与
\textbf{LTNE 双温度（局部非热平衡）}能量方程（\code{src/sub\_porous.f90}）：

\begin{align}
\nabla\!\cdot\!\mathbf{u} &= 0,\\
\nabla\!\cdot\!(\mathbf{u}\mathbf{u}) &= -\frac{\varepsilon^2}{\rho}\nabla p
  + \nu_{eff}\nabla^2\mathbf{u}
  - \frac{\varepsilon^2\mu}{K\rho}\mathbf{u}
  - \frac{\varepsilon^3 F}{\sqrt{K}}\,\lvert\mathbf{u}\rvert\,\mathbf{u},\\
\rho C_{p}\nabla\!\cdot\!(\mathbf{u}T_f) &= \varepsilon k_f\nabla^2 T_f + h_v\,(T_s-T_f),\\
0 &= (1-\varepsilon)k_s\nabla^2 T_s + h_v\,(T_f-T_s),
\end{align}
其中渗透率与 Forchheimer 系数按 \textbf{Ergun} 关系由颗粒尺度 $d_p$ 与孔隙率 $\varepsilon$ 计算：
\begin{equation}
K=\frac{d_p^{2}\varepsilon^{3}}{150(1-\varepsilon)^{2}},\qquad
F=\frac{1.75}{\sqrt{150}}\,\varepsilon^{1.5}.
\end{equation}
$h_v$ 为体积换热系数 [W/(m$^3$K)]；$h_v=0$ 时退化为局部热平衡（LTE，$T_s\equiv T_f$）。
变量约定与低速块一致：\code{U(1)=$\rho$}、\code{U(2..4)=速度}、\code{U(5)=T\_f}、压力在 \code{B\%p}，
骨架温度在 \code{B\%Ts}（初值 \code{Porous\_T\_ref}）。全部为 \textbf{SI}。

\subsection{多孔参数表}
\begin{table}[htbp]\centering\small
\caption{组 \code{\$porous\_ec} 与相关文件}\label{tab:porousparams}
\begin{tabular}{llp{8.0cm}}
\toprule
参数 & 默认值 & 说明/建议 \\
\midrule
\code{Porous\_T\_ref} & 288.15 & 骨架温度初值 [K]（也是 SIMPLE 路径的参考温度）\\
\code{Porous\_alpha\_Ts} & 0.7 & 骨架温度方程的欠松弛 \\
\code{Porous\_Max\_Iter} & 5000 & 每次调用多孔求解器的 SIMPLE/AC 迭代上限 \\
\code{Porous\_Tol} & 1e-8 & 压力修正残差判据 \\
\code{Porous\_U\_in/V\_in/W\_in} & 0,0,0 & 专用冷却剂入口速度 [m/s]；三者全 0 时\textbf{回落到 \code{LS\_U\_in/LS\_V\_in/LS\_W\_in}} \\
\code{Porous\_T\_in} & 0 & 冷却剂入口温度 [K]；$\le 0$ 时回落 \code{LS\_T\_ref} \\
\code{porous.inp} & --- & 逐块 $\varepsilon$、$d_p$[m]、$h_v$[W/(m$^3$K)]（见 \S 2.7）\\
\code{material.in} & --- & 类型码 3 + 骨架 $\rho,C_p,k_s$（SI）\\
\code{solid\_bc.inp} & --- & 骨架壁面热边界（等温 $T_w$ / 热流 $Q_w$）\\
\code{LS\_rho,mu,k,Cp} & 见表 \ref{tab:lsparams} & 流体物性直接复用低速参数 \\
\bottomrule
\end{tabular}
\end{table}

\paragraph{入口的优先级}
多孔块入口：若 \code{Porous\_U/V/W\_in} 至少一个非零 $\Rightarrow$ 用它作为入口速度；
否则用 \code{LS\_U\_in/LS\_V\_in/LS\_W\_in}。入口面仍需在 \code{bc3d.inp} 里标
码 5 / 21（入口），出口标 6 / 22。

\subsection{求解算法与预算}
\begin{itemize}
  \item \code{LS\_Algorithm=1/2}（SIMPLE/SIMPLEC）：\code{src/sub\_porous.f90}，含 DBF 源项；
        骨架温度与流体温度交替求解，\code{Porous\_alpha\_Ts} 控松弛。
  \item \code{LS\_Algorithm=3}（AC-FV，推荐用于耦合算例）：\code{src/sub\_porous\_ac.f90}
        与低速 AC 共用引擎，动量方程加入
        $c_{drag}=\bigl[\varepsilon^2\mu/K+\varepsilon^3F\rho\lvert\mathbf{u}\rvert/\sqrt{K}\bigr]/\rho$
        （残差项 $-\,c_{drag}V u_m$，LU-SGS 对角 $+\,c_{drag}V$）；流动收敛后再解 LTNE 温度对。
  \item \textbf{耦合预算}：交错耦合时 "多孔段" 的迭代量
        $= \code{AC\_Max\_Iter}\times\code{Porous\_Chunk\_Iter}$
        （SIMPLE 路径则为 $\code{Porous\_Max\_Iter}\times\code{Porous\_Chunk\_Iter}$）。
        实测参考：$5000\times3\approx170$ s/轮，$20000\times10\approx65$ min/轮
        $\Rightarrow$ \textbf{改这两个参数前先估时间}。
\end{itemize}

\subsection{参数选取建议}
\begin{itemize}
  \item \textbf{$\varepsilon$}：颗粒床 0.36--0.45；金属泡沫/丝网 0.8--0.95；
        本仓库 case3 取 0.30（致密多孔壁）。
  \item \textbf{$d_p$}：用实际颗粒/孔径 [m]；$K\propto d_p^2$，$d_p$ 偏大会低估阻力。
  \item \textbf{$h_v$}：由换热面积密度估算 $h_v\approx a_s h_{conv}$（$a_s$ 比表面积 [m$^2$/m$^3$]）；
        金属泡沫典型 $10^{5}$--$10^{7}$ W/(m$^3$K)。$h_v$ 越大 LTNE 温差越小。
  \item \textbf{收敛}：压力修正残差与骨架温度残差都要看；$\varepsilon$ 小、$d_p$ 小（阻力大）
        时需减小伪时间步（\code{AC\_CFL}）或增大 \code{Porous\_alpha\_Ts} 的欠松弛。
  \item \textbf{验证算例}：\code{porous\_plug}（Darcy/Darcy--Brinkman 解析）、
        \code{porous\_ltne\_1d}（LTNE 一维，含 $h_v$ 标定）、
        \code{beavers\_joseph}（码 16 界面滑移解析解）、
        \code{porous\_fluid\_phaseB}（码 19 发汗壁 + 分段交错）。
  \item \textbf{已知缺口}：多孔速度入口存在约 28\% 通量缺口（流量锁定问题）；
        压力求解器为迭代型，后续建议换 BiCGStab/ILU（见第 \ref{sec:validation} 章）。
\end{itemize}

\noindent\textbf{依据文件}：\code{src/sub\_porous.f90}、\code{src/sub\_porous\_ac.f90}、
\code{cases/porous\_plug/*}、\code{cases/beavers\_joseph/*}、\code{docs/程序能力与算例考核总结.md}。
%===EOF===

\section{固体导热求解器（\code{BLOCK\_SOLID}）}
\label{sec:solid}

\subsection{模型与离散}
块类型 1，\textbf{稳态导热}（\code{src/sub\_multi\_region.f90: solid\_solver\_one\_block}）。
有限体积法，逐格心热平衡
\begin{equation}
k\sum_{faces}\frac{A_{f}}{d_{f}}\,(T_{nb}-T_{i}) = 0 ,
\end{equation}
其中 $d_f$ 取\textbf{相邻格心连线的欧氏距离在面法向的投影}（含非正交修正，$\cos\phi$ 截断到 $\ge 0.1$），
因此支持非均匀网格与非直角曲线网格（O 网格等）。固体温度存于 \code{B\%Ts}（K，SI），
初值为 \code{T\_inf}（可压参考温度）。

\subsection{GS/SOR 迭代参数（组 \code{\$solid\_ec}）}
\begin{table}[htbp]\centering\small
\caption{固体求解控制}\label{tab:solidparams}
\begin{tabular}{llp{8.4cm}}
\toprule
参数 & 默认值 & 说明/建议 \\
\midrule
\code{Solid\_GS\_Omega} & 1.7 & SOR 过松弛因子（$1<\omega<2$；网格病态时降到 1.0--1.3）\\
\code{Solid\_Max\_Iter} & 20000 & 单次调用最大扫掠数 \\
\code{Solid\_Min\_Iter} & 5 & 允许判据生效前的最少扫掠数（此后每 100 次检查一次残差）\\
\code{Solid\_Tol} & 1e-9 & 收敛判据：$\max\lvert\Delta T_s\rvert$（K）\\
\bottomrule
\end{tabular}
\end{table}
日志形如 \code{Solid solver Block 1 GS iterations: 1400 final res: 8.8e-10}：
打印的扫掠数达到 \code{Solid\_Max\_Iter}$+1$ 说明\textbf{未达容差}（需增大上限、调 $\omega$ 或检查边界）。

\subsection{热边界（\code{solid\_bc.inp}）}
逐面给定：\textbf{$T_w>0$ 等温壁（K）；$T_w<0$ 给定热流 $Q_w$ [W/m$^2$]}；可选追加
\code{htc}/\code{Tinf} 列（对流边界）。未列出的面默认绝热（零梯度）。

\subsection{与流体/低速的共轭耦合（码 11 / 13）}
耦合例程在每个交界面按\textbf{热流平衡}计算界面温度（\code{compute\_interface\_T}）：
\begin{equation}
T_i=\frac{k_f\,T_f/dx_f + k_s\,T_s/dx_s}{k_f/dx_f + k_s/dx_s},
\qquad
\text{鬼点：}\ T_{ghost}=2T_i-T_{interior},
\end{equation}
其中 $dx$ 为格心到界面的距离，$k_f$ 已换算到流体无量纲体系
（$k_{ref}=\mu_{ref}c_{p,ref}/Pr_L$，$\mu_{ref}=\rho_{ref}Ma\,a_{ref}\texttt{Lscale}/Re$），
$k_s$ 为固体 SI 值（\code{material.in} 第 3 列）。方向/指标映射由连接描述符 \code{L1..L3} 完成，
支持任意 $(face,face1)$ 组合（$6\times6$）。

\paragraph{两种 CHT 模式（\code{Iflag\_Couple\_WallFlux}）}
\begin{itemize}
  \item \code{0}（默认，耦合式）：界面 $T_w$ 与 $q_w$ 由两侧场即时热流平衡得到；
  \item \code{1}（等温 $T_w$ / 热流 $q_w$ 壁通量式，\code{wfmode}）：气体侧把界面当\textbf{等温壁} $T_w$，
        固体侧施加热流 $q_w$，适合分段交错里"先算气体 $q_w$、再算固体"的流程（\code{cases/fluid\_solid}）。
\end{itemize}

\subsection{成本与使用建议}
\begin{itemize}
  \item \textbf{单次调用成本}：默认 $20000/10^{-9}$ 时，case1 规模（$67\times101\times2$）实测约
        \textbf{1400 次扫掠/次}；在\textbf{每时间步}都求解固体的强耦合模式下约 15--20 s/步
        （因此切强耦合时会自动把预算降到 \code{Solid\_Max\_Iter}$\le20$、\code{Solid\_Tol}$\ge10^{-6}$，
        见 \S\ref{sec:restart}）。
  \item \textbf{分段交错模式}下固体每外层轮只解一次，成本可接受，是长时程 CHT 的推荐用法。
  \item \textbf{稳定性}：界面热流很小时可适当放宽 \code{Solid\_Tol}（如 $10^{-7}$）以省时间；
        固体内部温差大/网格畸变严重时降低 $\omega$。
  \item \textbf{验证算例}：\code{solid\_1d}（一维线性解析，误差 $\sim10^{-6}$ K）、
        \code{solid\_annulus}（环形径向导热对数解，径向误差 $\sim1.9\times10^{-3}$ K）、
        \code{bl\_cht}（低速--固体共轭，码 13）。
\end{itemize}

\noindent\textbf{依据文件}：\code{src/sub\_multi\_region.f90}（\code{solid\_solver\_one\_block} /
\code{compute\_interface\_T} / \code{set\_solid\_ghost\_Bc}）、\code{src/sub\_read\_parameter.f90}（\code{read\_solid\_bc}）。
%===EOF===

\section{跨区域耦合}
\label{sec:coupling}

\subsection{两种耦合模式}
由 \code{\$couple\_ec: Iflag\_Couple\_Scheme} 选择：

\begin{table}[htbp]\centering\small
\caption{耦合模式}\label{tab:couplemode}
\begin{tabular}{clp{9.0cm}}
\toprule
取值 & 模式 & 说明 \\
\midrule
0 & \textbf{逐时间步同时耦合}（tight/强耦合） & 每个时间步把所有块各推进一次，然后调用
    \code{couple\_fluid\_solid\_interfaces} 等例程刷新交界面鬼点；气/固同步前进。\\
1 & \textbf{分段交错耦合}（staggered） & 主程序改走 \code{run\_staggered\_multiregion}：
    先推进一侧若干步（chunk），再解另一侧，反复交替并用外层判据收口；\code{t\_end} 不起作用。\\
\bottomrule
\end{tabular}
\end{table}

\paragraph{交错模式的调度表（\code{src/sub\_multi\_region.f90}）}
第 $it$ 轮（$it=1,\dots,\code{Niter\_Couple\_Outer}$）的气体步数
\begin{equation}
n_{step}(it)=
\begin{cases}
\code{Kstep\_Couple\_Comp}, & it\le \code{Niter\_Couple\_Warm},\\[2pt]
\max\!\bigl(\code{Kstep\_Couple\_Min},\ \code{Kstep\_Couple\_Comp}/2^{\,it-\code{Niter\_Couple\_Warm}}\bigr), & \text{否则},
\end{cases}
\end{equation}
即“先粗稳、再细修”。整段总步数 $=\sum_{it} n_{step}(it)$。
例如 case1（\code{Comp=1000, Warm=2, Outer=12, Min=1}）$=1000,1000,500,250,125,62,31,15,7,3,1,1=\mathbf{2995}$ 步。

\paragraph{外层收敛判据}
\begin{itemize}
  \item 码 11/13/19（含固体/多孔）：$\max\lvert\Delta T_w\rvert<\code{Tol\_Couple\_Tw}$（K）；
  \item 码 12（高低速流-流）：还要求 $\max\lvert\Delta p_w\rvert<\code{Tol\_Couple\_p}$（Pa）与
        $\max\lvert\Delta u_w\rvert<\code{Tol\_Couple\_u}$（m/s）；
  \item 判据从第 2 轮起生效（第 1 轮仅作初始化）；满足即提前退出，否则跑满 \code{Niter\_Couple\_Outer} 轮并打印
        \code{(not fully converged)}——\textbf{这是正常结束，不是崩溃}。
\end{itemize}

\subsection{交错耦合参数表}
\begin{table}[htbp]\centering\small
\caption{组 \code{\$couple\_ec}}\label{tab:coupleparams}
\begin{tabular}{llp{8.2cm}}
\toprule
参数 & 默认值 & 说明/选取建议 \\
\midrule
\code{Iflag\_Couple\_Scheme} & 0 & 0 同时耦合；1 分段交错 \\
\code{Kstep\_Couple\_Comp} & 1000 & 每轮气体段步数基准（\textbf{决定物理时间推进量}）\\
\code{Niter\_Couple\_Outer} & 60 & 外层轮数（\textbf{交错模式的时间推进总控制}）；太小则物理时间不足 \\
\code{Niter\_Couple\_Warm} & 0 & 前几轮跑满 \code{Kstep\_Couple\_Comp}（暖机）\\
\code{Kstep\_Couple\_Min} & 1 & 步数下限；\textbf{不要长期用 1}（末段几十轮只有 1 步，几乎不推进时间），
      建议 20--100 \\
\code{Porous\_Chunk\_Iter} & 10 & 每轮多孔段调用多孔求解器的次数（预算 $=\code{AC\_Max\_Iter}\times$ 本值）\\
\code{Twall\_Couple\_Init} & 300 & 首轮界面壁温初值 [K]（气体侧等温壁）\\
\code{Tol\_Couple\_Tw} & 2e-2 & 温度外层判据 [K]；与物理时间尺度匹配，否则永不满足 \\
\code{Tol\_Couple\_p}/\code{Tol\_Couple\_u} & 10 / 1e-2 & 码 12 的压力/速度判据 [Pa] / [m/s] \\
\code{Iflag\_Couple\_WallFlux} & 0 & CHT 模式：0 耦合式；1 等温 $T_w$/热流 $q_w$ 壁通量式（\code{wfmode}）\\
\code{Iflag\_Couple\_Restart} & 0 & 重启后按耦合状态续算（见 \S\ref{sec:restart}）\\
\bottomrule
\end{tabular}
\end{table}

\paragraph{交错模式的输出}
每轮结束写出 \code{iface\_couple.dat}（码 11/13/19：\code{x\_w, T\_w, q\_w, p\_w}）
或 \code{iface\_highlow.dat}（码 12：\code{x\_w, T\_w, p\_w, |u|\_w}），可直接画界面量收敛曲线。

\subsection{性能与选择指南}
\begin{table}[htbp]\centering\small
\caption{两种模式对比（\code{cases/fluid\_solid} 实测，单进程）}\label{tab:cost}
\begin{tabular}{p{3.4cm}p{5.2cm}p{5.2cm}}
\toprule
项目 & 逐时间步同时耦合（0） & 分段交错（1）\\
\midrule
固体求解频次 & \textbf{每步一次}（默认 $20000/10^{-9}$ 时约 15--20 s/步） & 每外层轮一次（$\sim$秒级）\\
气体步数控制 & \code{t\_end} / \code{Kstep\_save}（正常时间推进） & \code{Niter\_Couple\_Outer} 与调度表（\code{t\_end} 无效）\\
界面信息 & 每步即时更新（无滞后） & 外层轮间滞后（$n_{step}$ 步）；用判据控制滞后量\\
适用场景 & 界面变化快、强耦合、物理时间短 & 长时程 CHT/发汗冷却，界面缓变\\
强制要求 & --- & \textbf{必须单进程}（12/19 无跨进程耦合；码 11 有 \code{cht\_mpi\_fluid\_solid\_face}）\\
\bottomrule
\end{tabular}
\end{table}

\textbf{建议}：先用交错模式把界面"跑稳"（成本低），再视需求切到逐时间步强耦合；
切换与续算已由 \code{Iflag\_Couple\_Restart} 自动化（见 \S\ref{sec:restart} 与第 9 章 case 示例）。

\subsection{接口配对与连接描述符（新拓扑如何支持）}
\begin{itemize}
  \item 读入 \code{bc3d\_interface.inp} 后程序自动完成\textbf{块间配对}、\textbf{接口面指标对齐}
        （\code{src/sub\_convert\_inp.f90}、\code{src/sub\_IO.f90}）；
        链接式写法（物理码 + \code{nb1}/\code{face1}）会被"提升"为规范接口码，
        如 \code{bc<0} 或 $11\le bc\le 19$ 都按界面处理。
  \item 每个界面几何由描述符 $L_1,L_2,L_3$ 决定：本块第 $k$ 维 $\leftrightarrow$ 邻块第 $|L_k|$ 维，
        符号表示正/反向。\textbf{耦合例程一律用 L 表做指标映射}，
        因此像 case1 的 \code{gas i+ (face 4) <-> solid j- (face 2)} 这种"非同向"拓扑，
        只要该组合被允许即自动工作。
  \item 可压--固体通用耦合核 \code{couple\_compressible\_fluid\_solid\_face} 明确支持
        \textbf{任意 $6\times6$ 面组合}；码 11/12/19 的交错驱动在 2026-09-18 已扩展到
        \code{gas i+ <-> 对侧 j-}（\code{cases/fluid\_solid} 拓扑）。
  \item 自检：\code{IF\_Debug=1} 时启动会调用 \code{dbg\_dump\_interfaces}，
        打印每个界面的 \code{face/face1}、节点范围与 \code{L1..L3}——\textbf{新拓扑调试必用}。
\end{itemize}

\subsection{发汗壁（码 19 的 pair42 分支）}
当可压气体界面同时是\textbf{冷却剂注入面}时，启用发汗壁：
气体侧按\textbf{等温壁 $T_w$} $+$ 法向喷射通量处理（\code{wall\_bound\_blowing}），
多孔侧按界面温度外推骨架温度（$T_{s,g}=T_{s,i}+q_w dx/k_{s,eff}$）并返回新的 $T_w$。
控制参数：多孔入口由 \code{Porous\_U/V/W\_in}（或 \code{LS\_*} 回落）给定，
气体侧热流必须有粘性（\code{If\_viscous=1}）。

\subsection{常见问题（FAQ）}
\begin{description}
  \item[为什么只算了 2995 步就停了？] 交错模式的\textbf{设计终点}：跑满
    \code{Niter\_Couple\_Outer} 轮后正常退出（\code{t\_end} 无效）。要继续：
    ① 直接再跑一次（有 \code{field\_restart.dat} 会自动续算，界面量零跳变）；
    ② 增大 \code{Niter\_Couple\_Outer}、\code{Kstep\_Couple\_Min}、\code{Niter\_Couple\_Warm}。
  \item[外层判据永远不满足（判据平台高于 \code{Tol\_Couple\_Tw}）] 把判据放宽到平台之上
    （如 case1 平台 $\sim$0.1 K，可取 0.2），或接受"跑满轮数"的用法。
  \item[启动报 \code{File already opened in another unit}] 旧版 libgfortran 对“同一文件连两个 unit”
    会致命报错（新版放宽）。请使用 2026-09-25 之后修复的 \code{src/sub\_read\_parameter.f90}
    （该文件里 \code{scan\_control\_ec\_groups} 复用已打开的 unit，不再二次 \code{open}）。
  \item[\code{-np>1} 挂死/越界] 交错模式（码 12/19）必须 \code{mpirun -np 1}，
    且目录里不要放 \code{partation.dat}；码 11 支持跨进程 CHT。
  \item[NaN] 先查低速块：\code{cases/fluid\_solid} 里 \code{LS\_Algorithm=1}（SIMPLE）会立即发散，
    必须 AC-FV 且 \code{AC\_CFL=2}（详见 \S\ref{sec:lowspeed}）。
\end{description}

\noindent\textbf{依据文件}：\code{src/sub\_multi\_region.f90}、\code{src/sub\_convert\_inp.f90}、
\code{src/sub\_IO.f90}、\code{src/opencfd\_ec3d\_v1.16a.f90}、\code{memory-bank/constraints.md}。
%===EOF===


\section{输出、重启与耦合状态}
\label{sec:restart}

\subsection{输出文件总览}
\begin{table}[htbp]\centering\small
\caption{运行产物}\label{tab:outputs}
\begin{tabular}{p{4.1cm}p{3.6cm}p{7.0cm}}
\toprule
文件 & 触发 & 内容/单位 \\
\midrule
\code{Residual.dat} & \code{Kstep\_show} & 残差（追加写）\\
\code{force.log}、\code{force-invis-vis.log} & \code{Kstep\_show} & 气动力/力矩（粘性/无粘分开）\\
\code{flow3d.dat} & \code{Kstep\_save}（且 \code{Iflag\_savefile=1}） & 可压流场（PLOT3D 风格，无量纲）\\
\code{flow3d\_block\_*.vtk} / \code{flow3d.vtk} & \code{Kstep\_save} & VTK；\code{Iflag\_vtk\_onefile=0} 每块一个，$=1$ 合并；\code{Iflag\_vtk\_SI=1} 输出 SI \\
\code{flow3d\_node.dat} & \code{Kstep\_save} 且 \code{Iflag\_flow\_node=1} & \textbf{节点} SI 流场，6 变量 \code{d,u,v,w,T,Ts}，unformatted 逐块一条记录 \\
\code{Ts\_block\_*.dat} & \code{Kstep\_save} & 各固体/多孔块骨架温度（K）\\
\code{iface\_couple.dat} / \code{iface\_highlow.dat} & 交错每轮 & 界面量（$x_w,T_w,q_w,p_w$ / $|u|_w$）\\
\code{mesh-quality.dat}、\code{wall\_dist.dat} & 启动/自动 & 网格质量、壁距（缺失自动生成）\\
\code{field\_restart.dat} & \code{Kstep\_restart}（$\le0$ 时 $=\code{Kstep\_save}$） & \textbf{重启文件}（见 \S\ref{sec:restartfile}）\\
\code{output\_para.out} & 启动 & 全部生效参数 + 组来源回显（回归比对用）\\
\bottomrule
\end{tabular}
\end{table}

\subsection{重启文件 \code{field\_restart.dat}}
\label{sec:restartfile}
unformatted（双精度）布局：
\begin{lstlisting}
record 1 : iver, Total_block, Num_Mesh, NVAR, Kstep, LAP        (iver=2)
record 2 : tt, Ma, Re, gamma, T_inf, Lscale, dt_global
每块(全局块序) : nx,ny,nz
              : U (1:NVAR, 1-LAP:nx+LAP-1, ...)  含完整 LAP=4 鬼点缓冲
              : Ts, Tsn, p                      固体/骨架温度、上一步 Ts、低速/多孔压力
[iver=2 追加] trailer : Couple_State_Mode, Conv, Pair, Iter, nf, nk
                      : Couple_State_Twmax, Couple_State_Tol
                      : Tw_save, qw_save, u_save, pw_save   (nf x nk; 仅 nf,nk>0)
\end{lstlisting}
读入时校验 \code{nblock / NVAR / LAP / 块 1 维数 / Ma / Lscale}，不一致即\textbf{回退到原初始化}
（均匀场或 \code{Iflag\_init} 路径）；\code{Iflag\_restart=1}（强制）时不可用即 \code{stop 1}。
读入后自动刷新块间缓冲（\code{update\_buffer\_onemesh} / \code{update\_Ts\_buffer\_onemesh}），
因此高精度格式的缓冲信息被精确恢复。
\textbf{旧版本（\code{iver=1}，无 trailer）文件仍可读}，耦合状态视为"未知"（保持旧行为）；
被 kill 造成的残缺 trailer 也会被 \code{iostat} 容错回退。

\paragraph{开关}
\begin{itemize}
  \item \code{Iflag\_restart}：\code{0} 自动（存在即续算，默认）；\code{1} 强制（缺文件报错）；
        \code{-1} 完全关闭（不读也不写）；
  \item \code{Kstep\_restart}：写出间隔；$\le0$ 时等于 \code{Kstep\_save}；
  \item 注意重启文件较大（case1 约 17.5 MB/次），集群上请按需放大间隔。
\end{itemize}

\subsection{耦合状态与重启续算模式（\code{Iflag\_Couple\_Restart}）}
交错耦合的界面量（$T_w,q_w,u,p_w$）与收敛状态会随重启文件保存，于是续算可以：
\begin{itemize}
  \item \textbf{零跳变}：继续交错时\textbf{直接恢复界面量}，不再从 \code{control.ec} 初值重播；
  \item \textbf{自动切换}：若上次结束时已满足 \code{Tol\_Couple\_Tw}，则不再重播交错调度，
        而是把内部 \code{Iflag\_Couple\_Scheme} 置 0，进入\textbf{逐时间步强耦合}主循环
        （\code{t\_end} 重新生效），并在退出前补写一次重启文件。
\end{itemize}
\begin{table}[htbp]\centering\small
\caption{\code{Iflag\_Couple\_Restart}（组 \code{\$couple\_ec}，默认 0）}\label{tab:cprestart}
\begin{tabular}{clp{8.6cm}}
\toprule
值 & 行为 & 说明 \\
\midrule
0 & 自动 & 上次已满足判据 $\Rightarrow$ 切逐步强耦合；否则继续交错（界面量恢复）\\
1 & 强制强耦合 & 重启后立即进入逐时间步耦合（\code{t\_end} 生效）\\
2 & 强制交错 & 保持 2026-09-25 之前的行为（仍恢复界面量，避免跳变）\\
$-1$ & 忽略 & 不判定也不恢复界面量（\textbf{严格复现旧行为}，用于回归测试）\\
\bottomrule
\end{tabular}
\end{table}
\textbf{切强耦合时自动限制固体预算}：\code{Solid\_Max\_Iter}$=\min$(现值, 20)、
\code{Solid\_Tol}$=\max$(现值, $10^{-6}$)——只改扫掠上限与容差，SOR 因子与物理不变；
实际生效值打印在日志与 \code{output\_para.out}。

\subsection{节点流场 \code{flow3d\_node.dat}（\code{Iflag\_flow\_node=1}）}
\begin{itemize}
  \item 8 格心平均 $\to$ \textbf{网格节点}，\textbf{SI 单位}，unformatted，
        \textbf{逐块一条记录}，块序与 \code{Mesh3d.x} 一致，故可直接与 \code{Mesh3d.x} 组合显示；
  \item 每条记录 \textbf{6 个场}：\code{d, u, v, w, T, Ts}。
        \code{T} $=$ 流体温度（固体块为固体温度，沿用 \code{flow3d.vtk} 约定）；
        \code{Ts} $=$ \textbf{固体/骨架温度}（多孔块 LTNE 温差因此可见；流体/低速块镜像 \code{T}）；
  \item 校验工具：\code{python3 util/check\_flow3d\_node.py}（打印每块维度、点数与量级统计）。
\end{itemize}

\paragraph{MPI 约定（改代码时必守）}
读写文件\textbf{只由 0 号进程做}并遍历全局块序；其它 rank 只遍历\textbf{自己拥有}的块
（\code{B\_n(m)} 仅本进程有效，非 0 号进程访问全局块表会越界/挂死）。
新增开关必须同时进 \code{bcast\_para}（否则各 rank 取值不一致会在 I/O 处死锁）。

\noindent\textbf{依据文件}：\code{src/sub\_restart.f90}、\code{src/sub\_IO.f90}、
\code{docs/重启与节点流场输出说明.md}。
%===EOF===

\section{算例：\code{cases/fluid\_solid} 三个耦合算例}
\label{sec:cases}

三个算例共用同一套两分块网格，只改 \code{material.in}、\code{solid\_bc.inp} /
\code{porous.inp} 与 \code{control.ec}，分别在界面上演示\textbf{可压--固体（码 11）}、
\textbf{可压--低速（码 12）}、\textbf{可压--多孔（码 19）}。

\subsection{共同几何与拓扑}
\begin{itemize}
  \item \textbf{Block1}：$67\times101\times2$ 薄板（每块 $k$ 方向 2 层，即准二维）；
        $j=1$ 为对接面、$j=\max$ 为冷却端；$i\pm$ 为侧壁、$k$ 双向对称。
  \item \textbf{Block2}：$101\times201\times2$ 通道；$i=101$ 为对接面（分 3 段：壁 / 界面 / 壁），
        $i=1$ 为天花板壁，$j=1$ 为入口（码 5），$j=201$ 为出口（码 6），$k$ 双向对称。
  \item 界面共形：Block1 的 $j=1$ 面（67 个节点）与 Block2 的 $i=101$ 面 $j=68..134$ 段完全一致。
  \item 网格单位为 \textbf{mm}，故 \code{Lscale=0.001}。
  \item 运行时连接描述符（\code{IF\_Debug=1} 可打印，\code{docs/进度\_fluid\_solid三算例.md} 记录）：
\begin{lstlisting}
blk1 type1 bc=11 face=2 (j-)  nb1=2 face1=4 (i+)  L=(-2, 1, 3)
blk2 type0 bc=11 face=4 (i+)  nb1=1 face1=2 (j-)  L=( 2,-1, 3)
\end{lstlisting}
        即两侧物理法向都沿 $y$，但切向反向（$L_2=-1$）——
        “气体 $i+$” 对 “固体 $j-$” 的非同向拓扑，是 2026-09-18 专门扩展支持的组合。
  \item 因为法向都是 $y$，耦合公式里 \code{U(2)}=$x$ 切向、\code{U(3)}=$y$ 法向可与旧路径共用。
\end{itemize}

\subsection{共同的气体侧设置（\code{\$freestream\_ec} / \code{\$flow\_ec}）}
三个算例完全相同的部分：
\begin{lstlisting}
$freestream_ec
  Ma=3.0d0, Re=5000.d0, T_inf=800.d0, Lscale=0.001d0
$end
$flow_ec
  t_end=1.0d6, Kstep_save=500, Kstep_show=50,     ! t_end 在交错模式不使用
  dt_global=1.0d0, dtmax=1000.d0,
  Mesh_File_Format=2, IFLAG_LIMIT_FLOW=0, Bound_Scheme=-1,
  Iflag_vtk_onefile=1, Iflag_vtk_SI=1,
  Iflag_restart=0, Kstep_restart=0, Iflag_flow_node=1
$end
\end{lstlisting}
说明：\code{Twall=-1}（气体侧物理壁绝热，界面热流由耦合给出）、\code{If\_viscous=1}（默认，
热流必需）、\code{dt\_global=1.0} 使 \code{tt} 与 \code{Kstep} 同步计数（便于观察）。

\subsection{case1：可压--固体共轭传热（码 11）}
\paragraph{物理设置}
Block1 为不锈钢固体（\code{material.in}: \code{1 0} / \code{7900 500 16.3}），
\code{solid\_bc.inp} 指定其 $j+$ 面（面号 5）等温 300 K；Block2 为 $Ma=3$、$T_\infty=800$ K 空气。
热流从高温气体经界面导入固体、再由 $j+$ 面 300 K 冷端导出。

\paragraph{耦合设置}
\begin{lstlisting}
$couple_ec
  Iflag_Couple_Scheme=1, Niter_Couple_Warm=2, Niter_Couple_Outer=12,
  Tol_Couple_Tw=1.0d-2
$end
\end{lstlisting}
$\Rightarrow$ 气体段步数序列 $1000,1000,500,250,125,62,31,15,7,3,1,1$，
\textbf{一段共 2995 步}。外层判据 \code{Tol\_Couple\_Tw=1e-2} K。
固体 GS 用 \code{\$solid\_ec} 默认（$1.7/20000/5/10^{-9}$）。

\paragraph{预期日志}
\begin{lstlisting}
run_staggered_fluid_solid(paircode= 11): solid block 1  fluid block 2 ...
outer iter  1 : flow chunk steps = 1000
flow chunk done, Kstep= 1000  tt= 1000.0
solid chunk done, outer iter 1  max|dT_w|= 499.9 K  T_w range [799.90, 799.90]  max|q_w|= ...
...
Staggered CHT coupling reached Niter_Couple_Outer= 12 (not fully converged)
write field_restart.dat OK  (Kstep= 2995 , tt= 2995.0)
\end{lstlisting}
\textbf{注意}：本算例 \code{max|dT\_w|} 很快进入约 $0.0995$ K 的平台（固体尚未达稳态），
而判据是 $10^{-2}$ K，因此会打印 \code{(not fully converged)} 并跑满 12 轮结束——
这是\textbf{设定问题而非程序故障}（要早停需把判据放宽到 $\sim0.2$ K）。

\subsection{case2：可压--低速混接（码 12）}
\paragraph{物理设置}
Block1 改为\textbf{低速不可压空气}（\code{material.in}: \code{2 0}；低速物性由 \code{LS\_*} 给出），
$j=\max$ 面为冷却剂入口：目标质量通量 $G=2$ kg/(m$^2$s)、300 K。
Block2 仍为 $Ma=3$、$T_\infty=800$ K 空气。界面码 12。

\paragraph{关键换算（质量入口 $\to$ 等效速度入口）}
低速求解器为常密度（\code{LS\_rho=1.177} kg/m$^3$），且当前实现中 \code{LS\_Inlet\_Type=2}
按 $i$ 面处理，故本算例用\textbf{等效法向速度入口}：
\begin{equation}
\code{LS\_V\_in}=\frac{G}{\code{LS\_rho}}=\frac{2}{1.177}=1.699\ \text{m/s}\quad(+y).
\end{equation}

\paragraph{低速与 AC 设置（本算例的\textbf{强制}标定）}
\begin{lstlisting}
$lowspeed_ec
  LS_rho=1.177d0, LS_mu=1.846d-5, LS_k=0.0262d0, LS_T_ref=300.d0,
  LS_U_in=0.d0, LS_V_in=1.699d0, LS_P_in=0.d0, LS_P_out=0.d0,
  LS_alpha_p=0.2d0, LS_alpha_u=0.5d0,
  LS_Max_Iter=150, LS_Tol=1.d-7, LS_Algorithm=3
$end
$ac_ec
  AC_Max_Iter=5000
$end
$couple_ec
  Iflag_Couple_Scheme=1, Niter_Couple_Warm=2, Niter_Couple_Outer=12,
  Tol_Couple_Tw=2.0d0, Tol_Couple_p=1.0d2, Tol_Couple_u=1.0d-1
$end
\end{lstlisting}
\textbf{经验教训（务必保留）}：
\begin{itemize}
  \item \code{LS\_Algorithm=1}（SIMPLE）在本算例\textbf{第 1 步即发散}（\code{res\_p}$\sim10^{31}$），
        必须用 AC-FV（\code{=3}）；
  \item \code{AC\_CFL} 用默认附近的 \textbf{2}（\code{AC\_beta=10} 默认）；
        实测 \code{AC\_CFL=20, AC\_beta=100} 虽然 $res_q$ 更低，但\textbf{第 3 轮 NaN}；
  \item 码 12 的外层判据（$T_w$/$p_w$/$u_w$）比码 11 难满足，本算例整体放宽到
        $2.0$ K / $10^2$ Pa / $0.1$ m/s。
\end{itemize}

\paragraph{预期日志}
\begin{lstlisting}
run_staggered_highlow(12): gas block 2  LS block 1  Kstep_Couple_Comp= 1000 ...
outer iter 1 : gas chunk steps = 1000
highlow interface metrics, outer iter 1 : max|dT|= 321.0 K max|dp|= ... max|du|= ...
...
Staggered highlow coupling reached Niter_Couple_Outer= 12 (not fully converged)
\end{lstlisting}
界面量曲线见 \code{iface\_highlow.dat}。

\subsection{case3：可压--多孔（码 19，发汗冷却）}
\paragraph{物理设置}
Block1 改为\textbf{不锈钢多孔壁}（\code{material.in}: \code{3 0} + 骨架 \code{7900 500 16.3}），
\code{porous.inp} 给出块 1 的 $\varepsilon=0.30$、$d_p=10^{-3}$ m、$h_v=2\times10^{6}$ W/(m$^3$K)；
$j=\max$ 面同样是冷却剂入口 $G=2$ kg/(m$^2$s)、300 K（用 \code{Porous\_V\_in=1.699} m/s）。
界面码 19，采用\textbf{分段交错}：气体段按等温壁 $T_w$ 推进 $\to$ 提取 $q_w$ $\to$ 多孔段
（界面热流 $q_w$ + 底部注入）$\to$ 返回新的 $T_w$。

\paragraph{多孔与耦合设置}
\begin{lstlisting}
$porous_ec
  Porous_T_ref=300.d0, Porous_alpha_Ts=1.0d0, Porous_Max_Iter=30000,
  Porous_V_in=1.699d0, Porous_T_in=300.d0
$end
$ac_ec                       ! 多孔块用 AC 求解器
  AC_Max_Iter=5000, AC_Print=1000, AC_beta=100.d0, AC_CFL=20.d0, AC_Tol=1.d-6
$end
$couple_ec
  Iflag_Couple_Scheme=1, Niter_Couple_Warm=2, Niter_Couple_Outer=12,
  Porous_Chunk_Iter=3
$end
\end{lstlisting}
\textbf{预算压缩（实测）}：多孔段迭代量 $=\code{AC\_Max\_Iter}\times\code{Porous\_Chunk\_Iter}$
$=5000\times3\approx170$ s/轮；原先 $20000\times10$ 需 $\approx65$ min/轮且无必要。

\paragraph{预期日志}
\begin{lstlisting}
run_staggered_multiregion: block types F/L/S/P = T F F T
run_staggered_multiregion: pairs 11/12/13/19 = F F F T
gas block 2  porous block 1  pair42(gas i+/porous j-)= T
porous chunk done, outer iter 1  max|dT_w|= ... K  coolant exit |v|_max= ... m/s  T_w range [...]
\end{lstlisting}
界面量写 \code{iface\_couple.dat}；多孔块骨架温度另见 \code{Ts\_block\_1.dat} 与
\code{flow3d\_node.dat} 的 \code{Ts} 槽位（可看 LTNE 温差）。

\subsection{运行方式与时间预算（单进程实测）}
\begin{lstlisting}
cd cases/fluid_solid/run_case1_solid      # 或 run_case2_lowspeed / run_case3_porous
mpirun -np 1 ./opencfd-ec1.16a.out        # 交错耦合必须单进程
\end{lstlisting}
\begin{table}[htbp]\centering\small
\caption{三个算例的成本参考（本机单进程，含输出）}\label{tab:casecost}
\begin{tabular}{llp{8.2cm}}
\toprule
算例 & 单轮耗时 & 说明 \\
\midrule
case1 & $\approx100$ s & 气体段 1000 步 $\approx40$ s；固体每轮一次 GS（$\approx1400$ 扫掠）；全 12 轮 $\approx100$ s \\
case2 & $\approx105$ s & 由低速 AC 主导（\code{LS\_Max\_Iter=150} 每次调用）\\
case3 & $\approx170$ s & 由多孔 AC 主导（$5000\times3$）；$20000\times10$ 时约 65 min/轮 \\
\bottomrule
\end{tabular}
\end{table}
\textbf{输出体积}：\code{Iflag\_vtk\_onefile=1} 时每次 \code{Kstep\_save} 写合并 VTK $\approx$11 MB
（另加 \code{flow3d.dat} 3.3 MB、\code{field\_restart.dat} $\approx$17.5 MB/次）⇒
集群上建议放大 \code{Kstep\_save} 或关闭单文件 VTK。

\subsection{续算实测（重启耦合状态）}
以 case1 为例（2026-09-25 验证，\code{mpirun -np 1}）：
\begin{lstlisting}
# run1：跑满 12 轮，写 restart（Kstep=2995）
Staggered CHT coupling reached Niter_Couple_Outer= 12 (not fully converged)
write field_restart.dat OK  (Kstep= 2995 , tt= 2995.0)
# run2：自动续算（界面量零跳变，继续交错）
read_restart: coupling state  mode= 1 (0=tight,1=staggered)  satisfied= 0
              pair= 11  last max|dT_w|= 9.95e-2  interface cells= 66 x 1
restart: interface T_w/q_w restored from the restart file (no re-seeding; ...)
restart: continue staggered run at Kstep= 2995  tt= 2995.0
...
write field_restart.dat OK  (Kstep= 5990 ...)     # 每续一次再加 2995 步
\end{lstlisting}
若要\textbf{切换到逐时间步强耦合}继续算（\code{t\_end} 重新生效），二选一：
\begin{itemize}
  \item 在 \code{\$couple\_ec} 中写 \code{Iflag\_Couple\_Restart=1}（强制），并给 \code{t\_end} 一个合适值；
  \item 或把 \code{Tol\_Couple\_Tw} 放宽到平台之上（如 0.2 K），让上一段运行记录"已满足"，
        之后\textbf{自动}切换（日志会打印 \code{-> switch to per-step (tight) coupling}）。
\end{itemize}
切换后每次运行结束都会补写重启文件，可反复续跑（实测 \code{Kstep}: 150$\to$170$\to$190$\to$210）。

\subsection{调参清单（按现象查）}
\begin{table}[htbp]\centering\small
\caption{常见现象与处置}\label{tab:casefaq}
\begin{tabular}{p{5.2cm}p{8.4cm}}
\toprule
现象 & 处置 \\
\midrule
只跑 12 轮（$\approx2995$ 步）就退出 & 交错模式的\textbf{设计终点}：增大 \code{Niter\_Couple\_Outer}、
    \code{Kstep\_Couple\_Min}（如 100）、\code{Niter\_Couple\_Warm}，或直接续跑 \\
外层判据一直 \code{(not fully converged)} & 判据高于物理平台；放宽 \code{Tol\_Couple\_Tw} 或接受跑满轮数 \\
第 1 步就 NaN（case2/3） & 低速块误用 SIMPLE：必须 \code{LS\_Algorithm=3} 且 \code{AC\_CFL=2} \\
第 3 轮 NaN（case2） & \code{AC\_CFL/AC\_beta} 过大：回到 \code{AC\_CFL=2, AC\_beta=10} \\
case3 一轮太慢 & 压缩 \code{AC\_Max\_Iter}$\times$\code{Porous\_Chunk\_Iter}（如 $5000\times3$）\\
\texttt{-np 2} 挂死 & 交错模式（12/19）必须单进程；删掉目录里的 \code{partation.dat} \\
启动即 \code{File already opened in another unit} & 用 2026-09-25 之后修复的 \code{sub\_read\_parameter.f90}（见 \S\ref{sec:coupling} FAQ）\\
节点文件与网格不匹配 & 用 \code{util/check\_flow3d\_node.py} 检查；块序须与 \code{Mesh3d.x} 一致 \\
\bottomrule
\end{tabular}
\end{table}

\noindent\textbf{依据文件}：\code{cases/fluid\_solid/*/control.ec}（各目录另有 \code{control.ec.legacy} 旧写法备份）、
\code{docs/进度\_fluid\_solid三算例.md}、\code{docs/重启与节点流场输出说明.md}。
%===EOF===



\section{构建、运行与并行}
\label{sec:build}

\subsection{编译}
\begin{lstlisting}
cd src
make                 # 默认编译器 mpif90；产物 opencfd-ec1.16a.out
make clean           # 清理 *.o *.mod *.out
\end{lstlisting}
当前 \code{makefile} 的编译选项：
\begin{lstlisting}
f77 = mpif90
opt = -g -std=legacy -ffree-line-length-none -freal-4-real-8 -O3
\end{lstlisting}
其中 \code{-freal-4-real-8} 把单精度提升为双精度（\code{real(PRE\_EC)}），
\code{-std=legacy} 兼容旧代码风格。需要 OpenMP 时用 \code{-openmp} 并把 \code{NUM\_THREADS} 设为线程数。
\textbf{改动任何 \code{src/*.f90} 后必须重新 \code{make} 再验证。}

\subsection{运行}
\begin{lstlisting}
cd cases/<case>
mpirun -np 1 ./opencfd-ec1.16a.out          # 交错耦合（码 12/19）必须单进程
mpirun -np 4 ./opencfd-ec1.16a.out          # 多进程：需块数 >= 进程数，自动分区
\end{lstlisting}
\begin{itemize}
  \item 启动会依次做：读 \code{control.ec} $\to$ \code{prepare\_bc3d\_input}（校验/生成 \code{bc3d.inp}）
        $\to$ 网格质量检查 $\to$ \code{Init} $\to$ 读/生成 \code{wall\_dist.dat}
        $\to$ \code{Init\_flow}（有 \code{field\_restart.dat} 则续算）$\to$ 时间推进；
  \item 多进程时会自动生成 \code{partation-auto.dat}、\code{part\_grid.dat}；
        若目录里已有 \code{partation.dat} 则优先使用；
  \item \textbf{交错耦合模式目录里不要放 \code{partation.dat}}（否则可能被分到多进程而挂死）。
\end{itemize}

\paragraph{编译/运行环境注意事项（集群经验）}
\begin{itemize}
  \item \textbf{编译器的严格程度会暴露代码里违反标准之处}：例如“同一文件同时连到两个 unit”
        在新版 libgfortran 下被放宽（不报错），但旧版会致命报错
        \code{File already opened in another unit}。本仓库已在 2026-09-25 修复
        \code{control.ec} 的该问题；换工具链后若报类似信息，优先按此思路排查
        （\code{open} 同一文件两次/两个 unit）。
  \item \textbf{集群上通常是另一份源码副本}：改完要\textbf{同步 \code{src/} 并重新 \code{make}}，
        再把 \code{opencfd-ec1.16a.out} 拷到算例目录；不要只拷二进制而忘了源码版本一致。
  \item \textbf{保留 \code{field\_restart.dat}}：作业被时限杀死后可从最近保存点续算；
        建议把 \code{Kstep\_restart} 设小一些（如 100$\sim$500）以限制重算量。
  \item 作业脚本示例（PBS/Slurm 通用骨架）：
\begin{lstlisting}
cd $CASE_DIR
module load mpi/...          # 与编译时一致的 MPI
mpirun -np 1 ./opencfd-ec1.16a.out > run.log 2>&1
\end{lstlisting}
\end{itemize}

\subsection{算例目录产物与清理}
运行产物（可随时删除/重建）：\code{*.vtk}、\code{flow3d.dat}、\code{Residual.dat}、
\code{force*.log}、\code{*.inc}、\code{mesh-quality.dat}、\code{field\_restart.dat}、
\code{flow3d\_node.dat}、\code{Ts\_block\_*.dat}、\code{iface\_*.dat}、\code{partation-auto.dat}、
\code{part\_grid.dat}。
建议保留的“证据文件”：\code{run*.log}、\code{output\_para.out}、坏例的 \code{Residual.dat}。
仓库 \code{.gitignore} 已忽略 \code{field\_restart.dat} 与 \code{flow3d\_node.dat} 等大文件。

\subsection{调试手段}
\begin{itemize}
  \item \code{IF\_Debug=1}：打印接口配对诊断（\code{dbg\_dump\_interfaces}）；
  \item \code{Pdebug(1:4)}：分块/残差的额外打印级别；
  \item \code{Iflag\_bc\_check=2}：边界文件严格校验，便于定位 \code{bc3d.inp} 与
        \code{bc3d\_interface.inp} 不一致；
  \item \code{Kstep\_show}、\code{AC\_Print}：控制残差输出频率；
  \item 多进程时优先看 0 号进程日志；挂死常见原因是“各 rank 的开关取值不一致”
        （新开关必须进 \code{bcast\_para}）。
\end{itemize}

\noindent\textbf{依据文件}：\code{src/makefile}、\code{src/opencfd\_ec3d\_v1.16a.f90}、
\code{src/sub\_partation\_mpi.f90}、\code{memory-bank/constraints.md}。
%===EOF===

\section{验证、回归与已知问题}
\label{sec:validation}

\subsection{能力 $\leftrightarrow$ 算例对照（验证状态）}
\begin{table}[htbp]\centering\small
\caption{主要能力与其考核算例}\label{tab:validation}
\begin{tabular}{p{4.6cm}p{4.4cm}p{4.6cm}}
\toprule
能力 & 考核算例 & 状态/判据 \\
\midrule
可压跨声速气动 + SA & \code{DLR-F4}、\code{M6-wing} & 经典 Cp/气动力（原发行例）\\
混合 FV/FD（SEC） & \code{M6-wing-SEC} & 同上 + 嵌入区 \\
低速不可压核心 & \code{lidcavity}（Re=100） & Ghia 等 (1982) 中心线剖面，RMS 0.018 \\
低速三类入口 & \code{channel/type\{1,2,3\}} & Poiseuille 解析（$\bar U,u_{max},dp/dx$）\\
固体稳态导热 & \code{solid\_1d}、\code{solid\_annulus} & 解析解，误差 $10^{-6}$ K / $1.9\times10^{-3}$ K \\
多孔 Darcy / LTNE & \code{porous\_plug}、\code{porous\_ltne\_1d} & Darcy--Brinkman 解析 + LTNE 温差 \\
低速--多孔界面（16） & \code{beavers\_joseph} & Beavers--Joseph (1967) 滑移解析解 \\
可压--固体共轭（11） & \code{fluid\_solid/run\_case1\_solid} & 交错推进、界面量收敛、重启往返 \\
低速--固体共轭（13） & \code{bl\_cht} & 共轭传热联调 \\
可压--低速（12） & \code{high\_low\_fluid*}、\code{fluid\_solid/run\_case2} & 多块高低速稳定性 \\
可压--多孔（19） & \code{porous\_fluid\_phaseB}、\code{fluid\_solid/run\_case3} & 发汗壁/热耦合（物理收敛待细网格）\\
重启与节点输出 & \code{solid\_1d}、三例 fluid\_solid & Kstep 往返一致；节点文件与网格一致 \\
\bottomrule
\end{tabular}
\end{table}

\subsection{回归测试方法（改代码后必做）}
\begin{enumerate}
  \item \textbf{编译}：\code{cd src \&\& make}，确认 \code{EXIT=0}；
  \item \textbf{参数等价性}：对至少一个算例跑短程，用
\begin{lstlisting}
diff <(grep -v "namelist groups found" old/output_para.out) \
     <(grep -v "namelist groups found" new/output_para.out)
\end{lstlisting}
        比对生效参数；结构改动（如分组）应“除新增回显行外逐行一致”；
  \item \textbf{数值/稳定性}：\code{grep -i nan run.log} 应为 0；残差应下降；\code{EXIT=0}；
  \item \textbf{重启往返}：跑一段（写出 \code{field\_restart.dat}）后重启续算，
        确认 \code{Kstep/tt} 正确恢复、无 NaN、界面量零跳变；
  \item \textbf{节点输出}：\code{python3 util/check\_flow3d\_node.py} 核对块维度/点数与 \code{Mesh3d.x} 一致；
  \item \textbf{接口诊断}：新拓扑先用 \code{IF\_Debug=1} 检查 \code{L1..L3} 与 \code{face/face1}。
\end{enumerate}

\subsection{本手册涉及功能的验证记录（2026-09-25）}
\begin{itemize}
  \item \textbf{\code{control.ec} 双 unit 修复}：\code{make} EXIT=0；全部 \code{src} 内
        \code{control.ec} 仅剩 1 处 \code{open}；case1 整段 12 轮、\code{output\_para.out}
        与修改前逐字节一致；旧式单组算例（\code{solid\_1d}）正常。
  \item \textbf{重启耦合状态 + 强弱耦合切换}：\code{ct1}（12 轮等价 + 续算零跳变，Kstep 2995$\to$5990）、
        \code{ct2}（收敛后自动切逐步强耦合，150$\to$170 并补写重启）、\code{ct3a--d}
        （自动/强制强耦合/强制交错/忽略四种取值）全部 \code{EXIT=0}、\code{NaN=0}；
        \code{output\_para.out} 仅多 1 行新回显。
\end{itemize}

\subsection{已知问题与待办（摘要）}
\begin{itemize}
  \item \textbf{多孔速度入口}存在约 28\% 通量缺口（“流量锁定”未实现为按面法向的通量约束）；
  \item \textbf{低速压力求解器}为迭代型（Jacobi/GS），建议换 BiCGStab/ILU 以加速并提高鲁棒性；
  \item \textbf{低速网格病态}：\code{flatplate\_re1e6}、\code{high\_low\_fluid} 收敛慢/精度受损；
  \item \textbf{耦合外迭代}可加 Aitken/Anderson 加速；
  \item \textbf{接口 17（固--多孔）/18（多孔--多孔）未实现}；码 12/19 无跨进程耦合；
  \item \textbf{合并 VTK} 改分区输出以降低内存/IO 峰值；
  \item 三例 \code{cases/fluid\_solid/*/README.md} 待补充；
  \item \code{util/readflow3d-ver2.4a/2.5.f90} 仍存在“同一 \code{control.ec} 连两个 unit”的隐患
        （旧编译器会报错，修法与主程序一致）。
\end{itemize}

\noindent\textbf{依据文件}：\code{docs/程序能力与算例考核总结.md}（§3/§5/§6/§9）、
\code{docs/重启与节点流场输出说明.md}、\code{memory-bank/progress.md}。
%===EOF===

\appendix
\section{\code{control.ec} 参数总表}
\label{app:params}

本附录按组列出\textbf{全部}控制参数及其\textbf{代码默认值}（\code{src/sub\_read\_parameter.f90:
set\_default\_parameter}）。只要算例里不写，就用默认值；写错变量名会报错退出。
“无量纲”指按来流参考量归一化（\S 1.6）。

\subsection{\code{\$freestream\_ec}（24）}
\begin{longtable}{p{3.5cm}llp{6.6cm}}
\caption{来流/参考态/气体物性}\\
\toprule
参数 & 默认值 & 单位 & 说明 \\
\midrule
\endfirsthead
\toprule
参数 & 默认值 & 单位 & 说明 \\
\midrule
\endhead
\code{Ma} & 1.0 & --- & 来流马赫数 \\
\code{Re} & 1000 & --- & 雷诺数，参考长度 $=$ \code{Lscale} \\
\code{gamma} & 1.4 & --- & 比热比 \\
\code{AoA} & 0.0 & deg & 攻角 \\
\code{AoS} & 0.0 & deg & 侧滑角 \\
\code{T\_inf} & 288.15 & K & 来流温度（粘性参考、固体 \code{Ts} 初值）\\
\code{Twall} & $-1$ & K & 可压物理壁温；$<0$ 绝热 \\
\code{p\_outlet} & $-1$ & 无量纲 & 出口压力；$<0$ 外插 \\
\code{PrL} & 0.7 & --- & 层流 Prandtl 数 \\
\code{PrT} & 0.9 & --- & 湍流 Prandtl 数 \\
\code{Lscale} & 1.0 & m/grid & 长度尺度（mm 网格取 0.001）\\
\code{Ref\_S}/\code{Ref\_L} & 1.0 / 1.0 & --- & 气动力参考面积/长度 \\
\code{Centroid(1:3)} & 0,0,0 & --- & 力矩参考点 \\
\code{Cood\_Y\_UP} & 1 & --- & $y$ 轴朝上标志 \\
\code{Kt\_inf}/\code{Wt\_inf} & 1e-5 / 0.01 & --- & SST 初值 \\
\code{IF\_TurboMachinary} & 0 & --- & 涡轮机内流模式开关 \\
\code{Ref\_medium\_usrdef} & 0 & --- & 是否用户自定义参考介质 \\
\code{Turbo\_P0}/\code{Turbo\_T0} & 101330 / 288.15 & Pa / K & 涡轮机总压/总温 \\
\code{Turbo\_L0}/\code{Turbo\_w} & 1.0 / 0 & m / --- & 参考长度/转速 \\
\code{Turbo\_Periodic\_seta} & 0 & deg & 周期扇区角 \\
\bottomrule
\end{longtable}

\subsection{\code{\$flow\_ec}（52）}
\begin{longtable}{p{3.7cm}llp{6.3cm}}
\caption{流场控制：时间步/格式/模型/网格/输出/重启}\\
\toprule
参数 & 默认值 & 单位 & 说明 \\
\midrule
\endfirsthead
\toprule
参数 & 默认值 & 单位 & 说明 \\
\midrule
\endhead
\code{t\_end} & 100 & 无量纲 & 终止时间（\textbf{交错耦合模式下无效}）\\
\code{Kstep\_save} & 1000 & --- & 输出间隔（流场/VTK/节点/重启同节奏）\\
\code{Kstep\_show} & 1 & --- & 残差与气动力打印间隔 \\
\code{Kstep\_average} & 0 & --- & 时间平均输出间隔（0 关闭）\\
\code{Kstep\_smooth} & $-1$ & --- & 流场光滑间隔（$<0$ 关闭）\\
\code{Kstep\_init\_smooth} & 0 & --- & 初始阶段光滑间隔 \\
\code{CFL} & 1.0 & --- & 局部时间步 CFL \\
\code{dt\_global} & 0.01 & 无量纲 & 全局时间步 \\
\code{dtmax}/\code{dtmin} & 10 / 1e-9 & --- & 时间步上下限 \\
\code{Iflag\_local\_dt} & 1 & --- & 1 局部时间步；非定常请置 0 \\
\code{Time\_Method} & 0 & --- & 0 LU-SGS；1 Euler；3 RK3；$-1$ 双时间步 \\
\code{Step\_Inner\_Limit} & 20 & --- & 双时间步内迭代上限 \\
\code{Res\_Inner\_Limit} & 1e-10 & --- & 内残差判据 \\
\code{w\_LU} & 1.0 & --- & LU-SGS 松弛因子 \\
\code{If\_Residual\_smoothing} & 0 & --- & 残差光滑开关 \\
\code{If\_dtime\_mesh} & 1 & --- & 网格差时自动减小时间步 \\
\code{Iflag\_Scheme} & 5 & --- & 空间格式（0..9，见 \S 3.3）\\
\code{Iflag\_Flux} & 5 & --- & 无粘通量（1..6，见 \S 3.3）\\
\code{IFlag\_Reconstruction} & 0 & --- & 重构变量（0 原始/1 守恒/2 特征）\\
\code{Bound\_Scheme} & MUSCL2C & --- & 边界插值格式 \\
\code{IF\_Scheme\_Positivity} & 1 & --- & 正值保持开关 \\
\code{IFLAG\_LIMIT\_FLOW} & 1 & --- & 流场限制（防负密度/负压）\\
\code{Iflag\_turbulence\_model} & 0 & --- & 0 无/1 BL/2 SA/3 SST \\
\code{If\_viscous} & 1 & --- & 0 无粘/1 粘性（热流必需 1）\\
\code{Iflag\_init} & 0 & --- & 初场：0 均匀；1 读 \code{flow3d.dat}；$-1$ 置零 \\
\code{MUT\_MAX} & $-1$ & --- & 涡粘上限（$<0$ 不限）\\
\code{CP1\_NSA}/\code{CP2\_NSA} & 0.2 / 100 & --- & New-SA 常数 \\
\code{Ldmin,Ldmax} & 1e-6, 1000 & --- & 密度限制器 \\
\code{Lpmin,Lpmax} & 1e-6, 1000 & --- & 压力限制器 \\
\code{Lumax}/\code{LSAmax} & 1000 / 1000 & --- & 速度/湍流量限制 \\
\code{Mesh\_File\_Format} & 0 & --- & 0 unformatted；1 formatted；2 \code{Mesh3d.x} \\
\code{Num\_Mesh} & 1 & --- & 网格层数（多网格）\\
\code{Pre\_Step\_Mesh(1:3)} & 0 & --- & 各层预推进步数 \\
\code{NUM\_THREADS} & 1 & --- & OpenMP 线程数 \\
\code{IF\_Debug} & 0 & --- & 1 打印接口诊断 \\
\code{Pdebug(1:4)} & 1 & --- & 调试打印级别 \\
\code{Periodic\_dX/dY/dZ} & 0 & --- & 周期平移量 \\
\code{IF\_Innerflow} & 0 & --- & 内流模式开关 \\
\code{Iflag\_savefile} & 0 & --- & 1 写 \code{flow3d.dat} \\
\code{Iflag\_vtk\_onefile} & 0 & --- & 0 每块一个 VTK；1 合并 \code{flow3d.vtk} \\
\code{Iflag\_vtk\_SI} & 0 & --- & 1 VTK 输出 SI 单位 \\
\code{Iflag\_bc\_check} & 1 & --- & 0 旧行为；1 生成/校验；2 严格报错 \\
\code{Iflag\_restart} & 0 & --- & 0 自动续算；1 强制；$-1$ 关闭 \\
\code{Kstep\_restart} & 0 & --- & 重启写出间隔（$\le0$ 时 $=$ \code{Kstep\_save}）\\
\code{Iflag\_flow\_node} & 0 & --- & 1 写 \code{flow3d\_node.dat}（节点 SI 流场）\\
\bottomrule
\end{longtable}

\subsection{\code{\$lowspeed\_ec}（21，全部 SI）}
\begin{longtable}{p{3.5cm}llp{6.4cm}}
\caption{低速不可压块}\\
\toprule
参数 & 默认值 & 单位 & 说明 \\
\midrule
\endhead
\code{LS\_rho} & 1.2 & kg/m$^3$ & 常密度 \\
\code{LS\_mu} & 1.8e-5 & Pa$\cdot$s & 动力粘度 \\
\code{LS\_k} & 0.025 & W/(m$\cdot$K) & 导热系数 \\
\code{LS\_Cp} & 1005 & J/(kg$\cdot$K) & 比热 \\
\code{LS\_T\_ref} & 288.15 & K & 参考/初场温度 \\
\code{LS\_Inlet\_Type} & 1 & --- & 1 速度；2 质量；3 压力入口 \\
\code{LS\_U\_in}/\code{LS\_V\_in}/\code{LS\_W\_in} & 1 / 0 / 0 & m/s & 入口速度分量 \\
\code{LS\_Mdot\_in} & 0 & kg/s & 入口质量流量（类型 2）\\
\code{LS\_P\_in} & 101325 & Pa & 入口总压（类型 3）\\
\code{LS\_P\_out} & 101325 & Pa & 出口静压 \\
\code{LS\_T\_wall} & $-1$ & K & 壁温；$<0$ 绝热 \\
\code{LS\_U\_lid} & 0 & m/s & 顶盖驱动速度（$j+$ 面）\\
\code{LS\_alpha\_p} & 0.3 & --- & 压力欠松弛 \\
\code{LS\_alpha\_u} & 0.7 & --- & 速度欠松弛 \\
\code{LS\_alpha\_T} & 0.7 & --- & 温度欠松弛 \\
\code{LS\_Max\_Iter} & 5000 & --- & 每次调用内迭代上限 \\
\code{LS\_Tol} & 1e-8 & --- & 收敛判据（压力修正残差）\\
\code{LS\_Scheme} & 1 & --- & 1 一阶迎风/2 二阶迎风/3 MUSCL \\
\code{LS\_Algorithm} & 1 & --- & 1 SIMPLE/2 SIMPLEC/3 AC-FV \\
\bottomrule
\end{longtable}

\subsection{\code{\$ac\_ec}（15）}
\begin{longtable}{p{3.5cm}llp{6.4cm}}
\caption{人工压缩（AC）求解器}\\
\toprule
参数 & 默认值 & 单位 & 说明 \\
\midrule
\endhead
\code{AC\_Max\_Iter} & 40000 & --- & 伪时间迭代上限/次调用 \\
\code{AC\_Print} & 500 & --- & 打印间隔 \\
\code{AC\_beta} & 10 & --- & $\beta=\code{AC\_beta}U_{ref}^2$ \\
\code{AC\_CFL} & 2 & --- & 无粘 CFL（大值易发散）\\
\code{AC\_CFLv} & 0.5 & --- & 粘性 CFL \\
\code{AC\_Tol} & 1e-7 & --- & 收敛判据 \\
\code{AC\_w} & 1.0 & --- & LU-SGS 松弛 \\
\code{AC\_Flux} & 1 & --- & 1 Rusanov/2 Steger--Warming \\
\code{AC\_Recon} & 1 & --- & 1 MUSCL/2 WENO5/3 WENO3 \\
\code{AC\_Limiter} & 1 & --- & 1 van Leer/2 minmod \\
\code{AC\_WenoBlend} & 0 & --- & WENO 混入一阶通量比例 \\
\code{AC\_WallRecon} & 1 & --- & 壁面切向单侧二阶重构 \\
\code{AC\_WallP} & 1 & --- & 壁面压力二阶重构 \\
\code{AC\_MomDiss} & 0 & --- & 0 lumped/1 分裂混合/2 纯 $|u_n|$ \\
\code{AC\_MomFrac} & 0.2 & --- & 分裂混合保留声速比例 \\
\bottomrule
\end{longtable}

\subsection{\code{\$solid\_ec}（4）、\code{\$porous\_ec}（8）、\code{\$couple\_ec}（12）}
\begin{longtable}{p{4.1cm}llp{6.0cm}}
\caption{固体 / 多孔 / 耦合调度}\\
\toprule
参数 & 默认值 & 单位 & 说明 \\
\midrule
\endhead
\code{Solid\_GS\_Omega} & 1.7 & --- & 固体 SOR 因子 \\
\code{Solid\_Max\_Iter} & 20000 & --- & 固体最大扫掠数 \\
\code{Solid\_Min\_Iter} & 5 & --- & 判据生效前最少扫掠 \\
\code{Solid\_Tol} & 1e-9 & K & 固体收敛判据 $\max|\Delta T_s|$ \\
\code{Porous\_T\_ref} & 288.15 & K & 骨架温度初值 \\
\code{Porous\_alpha\_Ts} & 0.7 & --- & 骨架温度欠松弛 \\
\code{Porous\_Max\_Iter} & 5000 & --- & 多孔内迭代上限 \\
\code{Porous\_Tol} & 1e-8 & --- & 多孔收敛判据 \\
\code{Porous\_U\_in}/\code{V\_in}/\code{W\_in} & 0/0/0 & m/s & 多孔入口速度（全 0 时回落 \code{LS\_*}) \\
\code{Porous\_T\_in} & 0 & K & 多孔入口温度（$\le0$ 回落 \code{LS\_T\_ref}）\\
\code{Iflag\_Couple\_Scheme} & 0 & --- & 0 同时耦合/1 分段交错 \\
\code{Kstep\_Couple\_Comp} & 1000 & --- & 每轮气体步数基准 \\
\code{Niter\_Couple\_Outer} & 60 & --- & 外层轮数 \\
\code{Niter\_Couple\_Warm} & 0 & --- & 满步暖机轮数 \\
\code{Kstep\_Couple\_Min} & 1 & --- & 步数下限（建议 20$\sim$100）\\
\code{Porous\_Chunk\_Iter} & 10 & --- & 每轮多孔求解次数 \\
\code{Twall\_Couple\_Init} & 300 & K & 首轮界面壁温 \\
\code{Tol\_Couple\_Tw} & 2e-2 & K & 温度外层判据 \\
\code{Tol\_Couple\_p}/\code{u} & 10 / 1e-2 & Pa / m/s & 码 12 压力/速度判据 \\
\code{Iflag\_Couple\_WallFlux} & 0 & --- & 0 耦合式/1 等温 $T_w$+热流壁通量式 \\
\code{Iflag\_Couple\_Restart} & 0 & --- & 0 自动/1 强制强耦合/2 强制交错/$-1$ 忽略 \\
\bottomrule
\end{longtable}
%===EOF===


\section{编号、文件与单位速查}
\label{app:codes}

\subsection{面编号（全局统一）}
\begin{center}
\begin{tabular}{cccccc}
\toprule
1 & 2 & 3 & 4 & 5 & 6 \\
$i-$ & $j-$ & $k-$ & $i+$ & $j+$ & $k+$ \\
\bottomrule
\end{tabular}
\end{center}

\subsection{接口码}
\begin{center}\small
\begin{tabular}{clc}
\toprule
码 & 含义 & 状态 \\
\midrule
11 & 可压 $\leftrightarrow$ 固体（CHT） & ✅ \\
12 & 可压 $\leftrightarrow$ 低速 & ✅（无跨进程）\\
13 & 低速 $\leftrightarrow$ 固体 & ✅ \\
14 & 固固 & ✅（缓冲）\\
15 & 流流（同类） & ✅（缓冲）\\
16 & 低速 $\leftrightarrow$ 多孔 & ✅（已定量验证）\\
17 & 固 $\leftrightarrow$ 多孔 & ❌ 未实现 \\
18 & 多孔 $\leftrightarrow$ 多孔 & ❌ 未实现 \\
19 & 可压 $\leftrightarrow$ 多孔 & ✅（无跨进程）\\
\bottomrule
\end{tabular}
\end{center}
判定规则：\code{is\_interface\_bc(bc) = (bc < 0) .or. (11 <= bc <= 19)}。
方向不由接口码决定，而由几何配对与 \code{L1..L3} 决定。

\subsection{物理边界码}
\begin{center}\small
\begin{tabular}{cll}
\toprule
码 & 名称 & 说明 \\
\midrule
2 & \code{BC\_Wall} & 壁（可压 $T_{wall}<0$ 绝热；低速可用 \code{LS\_U\_lid}）\\
3 & \code{BC\_Symmetry} & 对称 \\
4 & \code{BC\_Farfield} & 远场（低速按出口处理）\\
5 & \code{BC\_Inflow} & 入口（低速按 \code{LS\_Inlet\_Type}）\\
6 & \code{BC\_Outflow} & 出口 \\
21 & \code{BC\_LS\_Inlet} & 低速入口（同 5）\\
22 & \code{BC\_LS\_Outlet} & 低速出口（同 6）\\
201 & \code{BC\_Wall\_Turbo} & 旋转壁 \\
401 & \code{BC\_Extrapolate} & 外插 \\
501 & \code{BC\_Periodic} & 周期（配 \code{Periodic\_d*}) \\
$-1$ & \code{BC\_Inner} & 块间内部界面 \\
$-2/-3$ & \code{BC\_PeriodicL/R} & 周期界面 \\
\bottomrule
\end{tabular}
\end{center}

\subsection{输入文件格式速查}
\begin{center}\small
\begin{tabular}{p{3.9cm}p{10.2cm}}
\toprule
文件 & 关键字段 \\
\midrule
\code{control.ec} & 7 组 namelist（旧组 \code{\$control\_ec} 兼容）；见附录 A \\
\code{bc3d.inp} & \code{nblock} / \code{ni nj nk} / 块名 / \code{nface} / \code{i1 i2 j1 j2 k1 k2 bc} \\
\code{bc3d\_interface.inp} & 同上 $+$ \code{nb1, face1, L1 L2 L3}（\code{bc=-1}；跨行续写用 $-1$ 占位）\\
\code{material.in} & \code{nblock} / 块类型行 / 每块 \code{rho Cp k} \\
\code{porous.inp} & \code{nblock} / \code{m eps dp[m] hv[W/(m$^3$K)]} \\
\code{solid\_bc.inp} & \code{nblocks} / 块号 / 面数 / \code{face Tw[K] Qw[W/m$^2$]} \\
\bottomrule
\end{tabular}
\end{center}

\subsection{单位与量纲速查}
\begin{center}\small
\begin{tabular}{lll}
\toprule
量 & 可压块 & 低速/多孔/固体块 \\
\midrule
长度（网格） & 网格单位（$x\cdot\code{Lscale}\to$ m） & 同左 \\
密度 & $\rho/\rho_\infty$ & kg/m$^3$（\code{LS\_rho}）\\
速度 & $u/U_\infty$ & m/s \\
温度 & $T/T_\infty$ & K \\
压力 & $p/(\rho_\infty U_\infty^2)$ & Pa（\code{B\%p}）\\
导热系数 & $\mu_{ref}c_{p,ref}/Pr_L$ 归一化 & W/(m$\cdot$K) \\
\bottomrule
\end{tabular}
\end{center}

\subsection{界面量公式（CHT，码 11/13）}
\begin{equation}
T_i=\frac{k_fT_f/dx_f+k_sT_s/dx_s}{k_f/dx_f+k_s/dx_s},\qquad
q_w=\frac{k_s\,(T_i-T_{s,i})}{dx_s}\ \ [\text{W/m}^2],
\qquad T_{ghost}=2T_i-T_{interior}.
\end{equation}
界面文件输出列：\code{iface\_couple.dat}（$x_w$[m], $T_w$[K], $q_w$[W/m$^2$], $p_w$[Pa]）、
\code{iface\_highlow.dat}（$x_w$, $T_w$, $p_w$, $|u|_w$）。

\subsection{日志关键字速查}
\begin{center}\small
\begin{tabular}{lp{9.4cm}}
\toprule
关键字 & 含义 \\
\midrule
\code{namelist groups found ...} & 各组来源（\code{T}=文件，\code{F}=默认）\\
\code{outer iter k : flow chunk steps = n} & 交错第 $k$ 轮气体步数 \\
\code{solid chunk done, outer iter ... max|dT\_w|=} & 界面温度变化（收敛判据量）\\
\code{... reached Niter\_Couple\_Outer= ... (not fully converged)} & 跑满轮数正常结束 \\
\code{restart from ... Kstep=} / \code{continue staggered run at Kstep=} & 续算起点 \\
\code{coupling state mode= ... satisfied= ...} & 重启文件里的耦合状态 \\
\code{switch to per-step (tight) coupling} & 自动/强制切强耦合 \\
\code{Solid solver Block ... GS iterations:} & 固体 GS 扫掠数（$=\code{Max\_Iter}+1$ 表示未收敛）\\
\code{write field\_restart.dat OK} & 重启文件写出 \\
\bottomrule
\end{tabular}
\end{center}
%===EOF===

\section{更新记录}
\label{app:changelog}

\paragraph{维护规则}程序每新增或修改功能，必须：
\begin{enumerate}
  \item 更新正文对应章节（含参数表）；
  \item 新增参数时同步附录 A；
  \item \textbf{在下表追加一条记录}（日期 / 版本 / 功能 / 涉及源码 / 验证算例与结果）；
  \item 重新编译本手册（\code{cd docs/程序说明 \&\& ./build.sh}）确认无错误。
\end{enumerate}

\begin{longtable}{p{2.0cm}p{1.3cm}p{4.3cm}p{3.3cm}p{3.4cm}}
\caption{版本与功能更新记录}\\
\toprule
日期 & 版本 & 功能/改动 & 涉及源码 & 验证 \\
\midrule
\endfirsthead
\toprule
日期 & 版本 & 功能/改动 & 涉及源码 & 验证 \\
\midrule
\endhead
2026-09-05$\sim$09-18 & 1.0 & 新增块类型 1/2/3（固体导热、低速不可压 SIMPLE/SIMPLEC、
    多孔 DBF+LTNE），跨区域接口码 11--19，CHT/发汗耦合例程，
    \code{cases/fluid\_solid} 三算例 & \code{sub\_multi\_region.f90}、
    \code{sub\_lowspeed.f90}、\code{sub\_porous.f90} 等 & \code{lidcavity}（Ghia）、
    \code{solid\_1d/annulus}（解析）、\code{beavers\_joseph}（BJ 解析）、码 11/12/19 联调 \\
2026-09-07 & 1.0 & 合并 VTK 输出 + SI 单位换算；低速/多孔变量约定统一
    （\code{U(2..4)}=速度） & \code{sub\_IO.f90}、\code{sub\_init.f90} & 算例重跑一致 \\
2026-09-08/17 & 1.0 & AC-FV 求解器（低速+多孔）：AC\_* 控制、$\beta$、\code{ac\_check\_controls}；
    移除不稳定的 AUSM+ 路径 & \code{sub\_lowspeed\_ac.f90}、\code{sub\_porous\_ac.f90} &
    \code{lidcavity\_ac}、\code{channel}、\code{beavers\_joseph/run\_ac} \\
2026-09-23 & 1.0 & \code{control.ec} 拆分为 7 组 namelist（旧组兼容、严格报错、来源回显）；
    新增 \code{\$solid\_ec} 四参数 & \code{sub\_read\_parameter.f90}、
    \code{sub\_modules.f90}、\code{sub\_multi\_region.f90} & 旧↔新 \code{output\_para.out}
    除新增回显外逐字节一致；两个 util 工具同步支持新格式 \\
2026-09-23 & 1.0 & 重启文件 \code{field\_restart.dat}（含 LAP=4 缓冲，\code{iver=1}）
    与节点 SI 流场 \code{flow3d\_node.dat}（6 变量） & \code{sub\_restart.f90}、
    \code{sub\_IO.f90}、\code{sub\_init.f90} & 单块/多进程续算往返；节点维度与 \code{Mesh3d.x} 一致 \\
2026-09-24 & 1.0 & \code{flow3d\_node.dat} 由 5 变量扩为 6 变量（补 \code{Ts} 骨架温度）；
    三例 \code{fluid\_solid} 迁移到 7 组写法 & \code{sub\_IO.f90}、
    \code{cases/fluid\_solid/*/control.ec} & \code{porous\_ltne\_1d} 可见 $\max|T-T_s|$；
    A/B 参数逐项一致 \\
2026-09-25 & 1.0.1 & \textbf{修复启动崩溃}：\code{scan\_control\_ec\_groups} 不再二次
    \code{open} \code{control.ec}（旧版 libgfortran 报 \code{File already opened in another unit}）&
    \code{sub\_read\_parameter.f90} & \code{make} EXIT=0；case1 整段 12 轮；
    \code{output\_para.out} 与修改前逐字节一致 \\
2026-09-25 & 1.1 & \textbf{重启耦合状态 + 强弱耦合切换}：\code{iver=2} trailer
    （mode/conv/pair/iter + 界面量）；\code{Iflag\_Couple\_Restart}（0 自动/1 强耦合/2 交错/$-1$ 忽略）；
    交错续算界面量零跳变；切强耦合时自动限制固体预算并补写重启 & \code{sub\_restart.f90}、
    \code{sub\_modules.f90}、\code{sub\_multi\_region.f90}、\code{opencfd\_ec3d\_v1.16a.f90}、
    \code{sub\_read\_parameter.f90} & \code{ct1}（12 轮等价 + 2995$\to$5990 零跳变）、
    \code{ct2}（自动切强耦合 150$\to$170）、\code{ct3a--d}（四种开关取值）；
    全部 EXIT=0 / NaN=0 \\
2026-09-25 & 1.1 & \textbf{本手册 v1.0 建立}（\code{docs/程序说明/}，LaTeX） & \code{docs/程序说明/*} &
    \code{xelatex} 编译通过 \\
\bottomrule
\end{longtable}

\paragraph{待写入本表的已知待办}
\begin{itemize}
  \item 多孔速度入口的“按面法向通量”实现（消除 $\sim$28\% 通量缺口）；
  \item 低速压力求解器改用 BiCGStab/ILU；
  \item 接口 17（固--多孔）/18（多孔--多孔）；码 12/19 的跨进程耦合；
  \item \code{util/readflow3d-ver2.4a/2.5.f90} 的同类“双 unit”隐患修复；
  \item 三例 \code{cases/fluid\_solid/*/README.md}。
\end{itemize}


\end{document}
